% ===========================================================================
%  Apendice A — Fontes Mermaid dos diagramas
% ===========================================================================
\chapter{Fontes Mermaid dos diagramas}
\label{ap:mermaid}

Todo diagrama deste documento é vetorial, desenhado nativamente em TikZ --- o que
garante tipografia única, nitidez em qualquer \emph{zoom} e ausência de
dependência externa de renderização. Este apêndice traz a fonte
\textbf{Mermaid equivalente} de cada figura, para quem precisa editar, revisar ou
versionar o diagrama fora do \LaTeX{} (num \emph{pull request}, num Storybook, num
quadro).

\begin{regra}
\textbf{Ao acrescentar um diagrama TikZ, inclua aqui a fonte Mermaid
equivalente}, referenciando o mesmo número de figura. Diagrama sem par Mermaid é
diagrama que só uma pessoa consegue alterar.
\end{regra}

\begin{nota}
As fontes abaixo usam \cd{C4Context}, \cd{C4Container} e \cd{C4Component} onde o
diagrama é C4, \cd{sequenceDiagram} onde é sequência, \cd{erDiagram} onde é
entidade-relacionamento, \cd{stateDiagram-v2} onde é máquina de estados,
\cd{classDiagram} onde é de classes (e também no perfil e na taxonomia da UML),
\cd{gantt} no diagrama de tempo e \cd{flowchart} nos demais. Não são idênticas ao
TikZ: o Mermaid não expressa posicionamento fino nem notas ancoradas, e as notas
longas dos diagramas TikZ aparecem aqui como \cd{note} ou comentário. Onde o
Mermaid não tem o tipo da UML (objetos, pacotes, componentes, estrutura composta,
perfil, comunicação, tempo, visão geral de interação e implantação), a primeira
linha da fonte é um comentário que diz a aproximação e o que ela perde.
\end{nota}

% ---------------------------------------------------------------------------
\section{Figura~\ref{fig:visao-geral} --- Visão geral da plataforma}

\begin{lstlisting}[language=mermaid]
C4Container
  title Visao geral da Plataforma Unooria

  Person(estudante, "Estudante", "1o EF ao 3o EM")
  Person(prof, "Professor", "")

  Container_Boundary(plat, "Plataforma Unooria") {
    Container(hub, "Hub", "Next.js", "Autentica, roteia por papel, embeda")
    Container(exp, "Experiencia", "Vite estatico + CDN", "Roda no iframe do navegador")
    Container(api, "API", "NestJS", "Regra, gabarito, tenancy")
    ContainerDb(pg, "PostgreSQL", "Cloud SQL", "RLS por escola")
    ContainerQueue(ps, "Pub/Sub", "", "Telemetria analitica")
    Container(age, "Agente do professor", "Python + ADK", "Vertex AI Agent Engine")
  }

  System_Ext(idp, "WSO2 Identity Server", "Identidade")
  System_Ext(vertex, "Vertex AI", "Gemini")

  Rel(estudante, hub, "HTTPS")
  Rel(prof, hub, "HTTPS")
  BiRel(hub, exp, "postMessage v1", "token, contexto")
  Rel_L(hub, idp, "OIDC")
  Rel(hub, api, "BFF")
  Rel(exp, api, "SDK: /telemetry, /ai, dominio")
  BiRel(api, pg, "SQL + RLS")
  Rel(api, ps, "publica")
  Rel(ps, api, "push POST /telemetry/ingest")
  Rel(api, vertex, "quota por perfil")
  Rel(api, age, "SSE")
  Rel(age, vertex, "Gemini, acervo")
\end{lstlisting}

% ---------------------------------------------------------------------------
\section{Figura~\ref{fig:casos-de-uso} --- Casos de uso}

\begin{lstlisting}[language=mermaid]
flowchart LR
  estudante([Estudante]); prof([Professor]); gestor([Gestor de escola]); admin([Gestor do contrato])
  idp[["Provedor de identidade (WSO2 IS de destino; hoje o hub)"]]
  vertex[["Vertex AI (Gemini e Agent Engine)"]]; pubsub[[Pub/Sub]]; ci[[CI/CD]]

  subgraph SIS["Plataforma Unooria"]
    UC01(("UC-01 Autenticar-se"))
    UC02(("UC-02 Escolher missao no hub"))
    UC03(("UC-03 Realizar experiencia"))
    UC04(("UC-04 Registrar telemetria"))
    UC05(("UC-05 Renovar token da sessao"))
    UC06(("UC-06 Ver o proprio resultado"))
    UC10(("UC-10 Preparar atividade"))
    UC11(("UC-11 Gerar conteudo com IA"))
    UC12(("UC-12 Conduzir o momento coletivo"))
    UC13(("UC-13 Ler resultados da turma"))
    UC14(("UC-14 Gerar plano de aula"))
    UC15(("UC-15 Check-in de bem-estar"))
    UC16(("UC-16 Zerar dados da turma"))
    UC45(("UC-45 Conversar com o agente de IA"))
    UC20(("UC-20 Importar pessoas e enturmar"))
    UC21(("UC-21 Editar pessoa e redefinir senha"))
    UC22(("UC-22 Ler agregado da rede"))
    UC23(("UC-23 Publicar manifesto"))
    UC24(("UC-24 Ativar ou reverter versao"))
  end

  %% UC-01 vale para todo papel; o desenho o liga so ao estudante.
  estudante --- UC01 & UC02 & UC03 & UC06
  prof  --- UC10 & UC14 & UC15 & UC45 & UC12 & UC13 & UC16
  gestor --- UC20 & UC21
  admin --- UC22 & UC23

  UC03 -.->|include| UC04
  UC03 -.->|include| UC05
  UC10 -.->|include| UC11
  UC24 -.->|extend| UC23

  UC01 --- idp
  UC11 --- vertex
  UC14 --- vertex
  UC15 --- vertex
  UC45 --- vertex
  UC04 --- pubsub
  UC23 --- ci
\end{lstlisting}

% ---------------------------------------------------------------------------
\section{Figura~\ref{fig:c4l1} --- C4 nível 1: contexto}

\begin{lstlisting}[language=mermaid]
C4Context
  title C4 nivel 1 - Contexto do sistema

  Person(estudante, "Estudante", "do 1o EF ao 3o EM")
  Person(prof, "Professor", "conduz a turma")
  Person(gestor, "Gestao da rede", "escola, regional e contrato")
  Person(admin, "Admin da Unooria", "opera a rede inteira")

  System(sis, "Plataforma Unooria", "Hub multi-tenant que autentica por papel, embeda experiencias imersivas versionadas, corrige e pontua no servidor, coleta telemetria, produz leitura de turma e da ao professor um agente de IA - com isolamento logico por escola")

  System_Ext(idp, "WSO2 Identity Server", "Identidade central, varios contratos")
  System_Ext(vertex, "Vertex AI", "Gemini, Agent Engine e Vertex AI Search")
  System_Ext(ci, "GitHub Actions", "Publica imagens, bundles e manifestos")
  System_Ext(tf, "Terraform (repo separado)", "Provisiona a stack do contrato")
  System_Ext(avatar, "LiveAvatar (HeyGen)", "Escuta, voz e video do atendente da Ajuda (ADR-132)")

  Rel(estudante, sis, "realiza missoes")
  Rel(prof, sis, "libera, conduz, le; conversa com o agente")
  Rel(gestor, sis, "administra pessoas; le o agregado")
  Rel(admin, sis, "liga recursos; tetos de IA")
  Rel(idp, sis, "OIDC")
  BiRel(sis, vertex, "HTTPS + ADC")
  Rel(ci, sis, "publica")
  Rel(tf, sis, "provisiona")
  BiRel(sis, avatar, "sala; /ajuda/llm")
  %% Mao dupla: a API pede a sala ao vendor e o vendor chama de volta o LLM da
  %% API. O navegador tambem fala direto com o vendor (video por WebRTC).
\end{lstlisting}

% ---------------------------------------------------------------------------
\section{Figura~\ref{fig:c4l2} --- C4 nível 2: contêineres}

\begin{lstlisting}[language=mermaid]
C4Container
  title C4 nivel 2 - Conteineres (stack de um contrato)

  Person(estudante, "Estudante")
  Person(prof, "Professor")
  Person(adm, "Gestao e admin")

  Container_Boundary(stack, "Plataforma Unooria - stack de um contrato") {
    Container(hub, "Hub", "Next.js App Router", "Sessao, roteamento por papel, chrome institucional, montagem do embed, BFF (inclusive o stream do chat do agente), JWKS e assinatura de token")
    Container(exp, "Experiencia", "Vite + React, ESTATICO", "Uma por missao. Roda no navegador; nunca e um processo do servidor")
    Component(sdk, "@edu/sdk", "TypeScript", "Ponte, telemetria, API, progresso")
    Container(api, "API", "NestJS + Prisma", "Autorizacao por papel, tenancy, regra de dominio, gabarito, correcao, gateway de IA, ingestao de telemetria")
    ContainerDb(pg, "PostgreSQL", "Cloud SQL HA + PITR", "Unico banco, unico schema; RLS por escola")
    ContainerQueue(ps, "Pub/Sub", "", "Canal analitico de telemetria")
    Container(buckets, "Object Storage", "Cloud Storage", "bundles, assets publicos, planos privados")
    Container(engine, "Agente do professor", "Python + ADK, Vertex AI Agent Engine", "Sessao, ferramentas, guarda de escopo e titulo; so a API o chama (ADR-126)")
    ContainerDb(ds, "Vertex AI Search", "dois data stores, global", "habilidades (estruturado) e acervo (PDFs) - ADR-127")
  }

  System_Ext(vertex, "Vertex AI", "Gemini e Gemini Live")
  Container_Ext(idp, "WSO2 IS", "", "Externo a stack")
  System_Ext(avatar, "LiveAvatar", "HeyGen; video da Ajuda (ADR-132)")

  Rel(estudante, hub, "HTTPS")
  Rel(prof, hub, "HTTPS")
  Rel(adm, hub, "HTTPS")
  BiRel(hub, exp, "postMessage v1 - targetOrigin explicito")
  Rel(exp, sdk, "usa")
  Rel(sdk, api, "HTTPS + Bearer")
  Rel(hub, api, "BFF")
  BiRel(api, pg, "SQL + RLS")
  Rel(api, ps, "publica")
  Rel(ps, api, "push /telemetry/ingest")
  Rel(api, buckets, "assets, planos")
  Rel(api, vertex, "HTTPS; WSS do Gemini Live (ADR-158)")
  Rel(api, engine, "SSE (:streamQuery), ADC")
  Rel(engine, vertex, "Gemini")
  Rel(engine, ds, "busca")
  Rel(hub, idp, "OIDC")
  Rel(buckets, exp, "CDN")
  BiRel(api, avatar, "sala; /ajuda/llm")
  %% O video vai do vendor direto ao navegador (WebRTC), sem passar pela stack;
  %% a chave da conta, o prompt e a base nunca saem da API.
\end{lstlisting}

% ---------------------------------------------------------------------------
\section{Figura~\ref{fig:c4l3api} --- C4 nível 3: componentes da API}

\begin{lstlisting}[language=mermaid]
C4Component
  title C4 nivel 3 - Componentes da API (NestJS)

  Container_Boundary(api, "API") {
    %% --- transversal: roda em toda rota
    Component(guard,  "OidcGuard", "", "Valida JWT por JWKS")
    Component(roles,  "RolesGuard", "", "Restringe por papel")
    Component(tenint, "TenancyInterceptor", "", "Preenche escola/perfil no AsyncLocalStorage aberto pelo middleware")
    Component(obs,    "observabilidade/", "pino + OpenTelemetry", "Contexto da requisicao (middleware), logger com redaction, filtro global, sampler por rota; otel.ts e a primeira importacao do main.ts")
    %% validarEnv (Zod no boot) nao esta na figura TikZ

    %% --- plataforma: serve a todas as experiencias
    Component(mais,   "analise/ telemetry/ progress/ portabilidade/ prospeccao/ health/", "", "Tambem de plataforma; sem setas na figura (prospeccao/ sem tenant, ADR-107)")
    Component(ident,  "auth/ usuarios/ rostering/", "", "Identidade: TokenService; pessoas, escolas, turmas")
    Component(expmod, "experiencias/", "", "Catalogo e versao")
    Component(recs,   "recursos/", "", "Disponibilidade (ADR-87), quotas de IA (ADR-105), limites e configuracao da Ajuda, custos de IA (ADR-161); escrita so admin")
    Component(salas,  "salas/", "", "Liberacao por turma (ADR-69)")
    Component(agente, "agente-professor/", "", "Agente de IA do professor: posse da conversa, SSE; @Produto('agente-professor') (ADR-126)")
    Component(ajuda,  "ajuda/", "", "Avatar (ajuda-llm, fora dos guards de pessoa) e mascote (WS /ajuda/mascote/ao-vivo); por minutos, fora do QuotaService (ADR-132, ADR-158)")
    Component(equipes,"equipes/", "", "Quem divide um computador na missao (ADR-111); consulta.ts e a porta")
    Component(evid,   "evidencias/", "", "Habilidade da BNCC por atividade e rodada (ADR-112); a porta marcar")
    Component(podcast,"podcast/", "", "Recapitulacao de toda missao (ADR-150); registrarFonte na subida; @Global")

    %% --- infraestrutura: as portas unicas
    Component(prisma, "PrismaService", "Prisma", "comTenancy: transacao + RLS")
    Component(comum,  "comum/", "", "Guardas de liberacao (exigirMissaoLiberada) e de disponibilidade; series, rate limit")
    Component(ai,     "ai/", "", "AiProvider medido (MedidorDeProvider), QuotaService, AgenteProfessorPort, AvatarAoVivoPort, VozAoVivoPort")
    Component(store,  "object-storage/ + assets/", "", "local | GCS")

    %% --- dominio por experiencia
    Component(missoes, "motor/ geracao/ dossie/ atividades/ perfis/ unoogpt/ a-selecao/ jornal-escola/ juri-estadual/ quem-mandou-isso/", "", "Missoes com ficha, equipe e podcast: exp08ef1b, exp02em1b, exp07ef1b, exp04ef1b, exp06ef1b, exp05ef1b")
    Component(outras,  "vila/ rastro-boato/ manifesto-ia/ foto-conta/", "", "Outras missoes: app-70, exp06ef2b, exp08ef4b, exp01ef1b; sem seta na figura")
    Component(apps,    "planos/ bemestar/ planejamento/", "", "Apps do professor: app-90, app-91, app-92")
  }

  %% --- relacoes desenhadas na figura TikZ
  Rel(guard, roles, "antes de")
  Rel(tenint, prisma, "escola do token")
  Rel(salas, equipes, "contagem na sala")
  Rel(evid, equipes, "gancho aoMudarEquipes + consulta.ts")
  Rel(agente, ai, "AgenteProfessorPort; QuotaService: temSaldo antes, consumir no finally")
  Rel(ajuda, ai, "AiProvider.conversarEmFluxo, AvatarAoVivoPort, VozAoVivoPort; FORA do QuotaService")
  Rel(podcast, ai, "generateSpeech + generateJson; QuotaService")
  Rel(missoes, equipes, "consulta.ts (painel por computador; perfis/ registra a fonte de lado)")
  Rel(missoes, evid, "marcar na transacao da jogada + registrarContexto")
  Rel(missoes, podcast, "registrarFonte na subida; gravarFalasDoPodcast na homologacao")
  Rel(missoes, ai, "IA + quota")
  Rel(apps, ai, "IA + quota")
  Rel(apps, store, "bucket privado (planos, planejamento)")

  %% --- relacoes fora da figura TikZ (de todo modulo com dado ou de toda missao)
  Rel(salas, prisma, "comTenancy")
  Rel(equipes, prisma, "comTenancy")
  Rel(evid, prisma, "comTenancy")
  Rel(agente, prisma, "comTenancy: agente_professor_conversas")
  Rel(ajuda, prisma, "ajuda_sessoes, ajuda_limites: fora de comTenancy, por usuario_id")
  Rel(ajuda, comum, "exigirProdutoDisponivel (a cada turno do LLM)")
  Rel(podcast, prisma, "comTenancy: podcast_roteiros")
  Rel(podcast, store, "trechos e episodio no bucket")
  Rel(missoes, prisma, "comTenancy")
  Rel(missoes, comum, "exigirMissaoLiberada")
  Rel(outras, prisma, "comTenancy")
  Rel(outras, comum, "exigirMissaoLiberada")
  Rel(outras, ai, "manifesto-ia/ e foto-conta/: IA + quota")
  Rel(apps, prisma, "comTenancy")
  Rel(recs, prisma, "disponibilidades, ia_quotas, ajuda_limites, ia_precos")
  Rel(missoes, outras, "perfis/ importa o VilaModule (POST /atividades/:id/perfis/da-vila): a unica dependencia de modulo entre dominios")

  %% A figura TikZ poe a infraestrutura numa COLUNA a direita (c4), separada
  %% por uma pista: plataforma e dominio chegam a ela na horizontal, sem uma
  %% faixa atravessar a outra.
\end{lstlisting}

% ---------------------------------------------------------------------------
\section{Figura~\ref{fig:c4l3hub} --- C4 nível 3: componentes do hub}

\begin{lstlisting}[language=mermaid]
C4Component
  title C4 nivel 3 - Componentes do Hub (Next.js)

  Container_Boundary(hub, "Hub") {
    %% --- borda
    Component(papeis,  "lib/papeis.ts", "", "edge-safe: papel -> area")
    Component(mw,      "middleware.ts", "", "Sessao + area do papel")
    Component(authlib, "lib/auth.ts", "NextAuth", "A sessao")

    %% --- rotas (app/)
    Component(publicas, "Rotas publicas", "", "/ (landing, ADR-40), /login, /privacidade, /termos; fora do matcher")
    Component(areas,    "Areas por papel", "", "/estudante, /professor, /admin, /analise, /unooria; o layout faz trilha, largura e escola ativa")
    Component(paginas,  "Paginas com embed", "", "/exp/:expId, /dev-embed, /professor/plano-de-aula, /professor/plano-de-ensino, /professor/bem-estar")
    Component(jwks,     "/.well-known/jwks.json", "", "A chave publica")
    Component(bff,      "/api/plataforma/[...path]", "", "BFF generico, teto de 180 s")
    Component(stream,   "/api/plataforma/agente-professor/conversas/[id]/mensagens", "", "Rota de stream do turno do agente (ADR-128)")

    %% --- lib/ no servidor
    Component(embed,   "lib/embed-experiencia.ts", "", "Token + manifesto + origem, UMA vez")
    Component(devsrv,  "lib/experiencias-dev.ts", "", "Porta do dev server")
    Component(chave,   "lib/chave-token.ts", "", "Assina RS256, TTL 15 min")
    Component(apiserv, "lib/api-servidor.ts", "", "A API vista dos Server Components")
    Component(bffts,   "lib/bff.ts", "", "Passagem da sessao e token curto, dividida pelas duas rotas")

    %% --- components/
    Component(casca,   "CascaPlataforma", "", "CSS do DS + fontes")
    Component(header,  "HeaderPlataforma", "", "Um so para todos os papeis")
    Component(exphost, "ExpHost", "", "DispositivoFrame + ExpFrame (sandbox + protocolo)")
    Component(ajuda,   "components/ajuda/", "", "Janela da Ajuda: avatar e mascote; a sala abre na montagem e fecha na desmontagem")
    Component(chat,    "components/agente-professor/", "", "Chat do agente, montado pelo layout de /professor")
  }
  Container_Ext(api, "API", "NestJS", "Valida o token pelo JWKS do hub")

  Rel(mw, papeis, "usa")
  Rel(mw, authlib, "le a sessao")
  Rel(paginas, embed, "montarEmbed")
  Rel(embed, exphost, "src, token e origem, pela pagina")
  Rel(embed, devsrv, "porta em dev")
  Rel(embed, chave, "assina")
  Rel(jwks, chave, "chave publica")
  Rel(apiserv, chave, "assina o token curto")
  Rel(bffts, chave, "token curto")
  Rel(header, ajuda, "item Ajuda do menu da conta: nao estudante e produto ligado")
  Rel(ajuda, bff, "abrir e encerrar a sessao")
  Rel(bff, bffts, "usa")
  Rel(stream, bffts, "usa")
  Rel(chat, stream, "o turno (SSE)")
  Rel(bff, api, "HTTPS + Bearer")
  Rel(stream, api, "SSE")
  Rel(apiserv, api, "HTTPS + Bearer")
  Rel(ajuda, api, "WSS do mascote, com o bilhete que o BFF entregou")

  %% Fora da figura TikZ: o chat lista e apaga conversas pelo BFF generico;
  %% CascaPlataforma e usada pelos layouts, por /exp e pela landing.
  Rel(chat, bff, "listar, criar, ler e apagar conversas")
\end{lstlisting}

% ---------------------------------------------------------------------------
\section{Figura~\ref{fig:c4l4} --- C4 nível 4: tipos que carregam a regra}

\begin{lstlisting}[language=mermaid]
classDiagram
  class connectToHub {
    +hubOrigin: string
    +timeoutMs = 8000
    +READY_RETRY_MS = 400
    valida ev.origin, senao descarta
    reenvia exp:ready ate responder
    compara protocolVersion MAIOR
    valida pos-init com Zod
  }

  class SdkTypes {
    +PROTOCOL_VERSION = 2
    +HubInitSchema
    +HubParaExpSchema
    +ExpParaHubSchema
    unioes discriminadas por type
  }

  class PrismaService {
    +comTenancy(fn, escolaId?, opts?)
    abre $transaction
    set_config('app.escola_id', ..., true)
    sem escola: LANCA
  }

  class TenancyContext {
    +AsyncLocalStorage~ContextoDaRequisicao~
    +comEscolaNoContexto(id, fn)
    +escolaDoContexto(): string?
  }

  class LinhaDeLeitura {
    <<union>>
    grau scored: numero: number
    grau indicative | zero | none: numero: null
  }

  connectToHub ..> SdkTypes : usa
  PrismaService ..> TenancyContext : le
  note for LinhaDeLeitura "Comentario avisa; TIPO impede.\nO hub nunca recebe numero onde nao ha medida."
\end{lstlisting}

% ---------------------------------------------------------------------------
\section{Figura~\ref{fig:monorepo} --- Grafo de dependências do \emph{monorepo}}

\begin{lstlisting}[language=mermaid]
flowchart LR
  ds[["@unooria/design-system (npm publico, repo storybook-unooria)"]]
  ui["packages/ui"]
  sdk["packages/sdk"]
  exps["Experiencias com interface: exp01ef1b, exp02em1b, exp04ef1b, exp05ef1b,
        exp06ef1b, exp06ef2b, exp07ef1b, exp08ef1b, exp08ef4b; app-90, app-91,
        app-92; _template e backup-exp04ef1b"]
  hub["apps/hub"]
  e70["app-70"]
  api["apps/api"]; agente["apps/agente-professor (Python + ADK)"]

  ds -->|versao publicada| ui
  ds -->|CSS e fontes| exps & hub
  ui --> exps & hub
  sdk --> exps & hub & e70
  %% apps/api e apps/agente-professor NAO importam pacote do workspace
  %% app-70 NAO depende de packages/ui: canvas 3D, sem DOM de UI (ADR-18)
\end{lstlisting}

% ---------------------------------------------------------------------------
\section{Figura~\ref{fig:seq-feed} --- Curadoria de um bloco de \emph{feed}}

\begin{lstlisting}[language=mermaid]
sequenceDiagram
  autonumber
  participant E as Experiencia (exp08ef1b no iframe)
  participant C as runs.controller (casca HTTP)
  participant S as motor/servico.ts (regra + persistencia)
  participant P as PostgreSQL

  Note over E,C: durante o feed
  E->>C: POST /runs/:id/eventos
  C->>P: grava os eventos novos (comTenancy, idempotente)
  C-->>E: gravados, descartados

  Note over E,C: a cada bloco
  E->>C: GET /runs/:id/feed
  C->>P: a run e do estudante? (comTenancy)
  C->>S: proximoBlocoDoFeed
  S->>P: comTenancy (escola)
  S->>P: pg_advisory_xact_lock
  Note over S,P: Sem o lock, dois pedidos simultaneos de bloco<br/>se sobrescrevem e um bloco inteiro desaparece.<br/>Transacao com timeout de 20s.
  P-->>S: run, eventos, catalogo
  S->>S: aplicarEventos e curarBloco (puras)
  S->>P: UPDATE engine_state
  S->>S: paraFeedPost: sem lado, perfil nem intensidade
  Note over S: Esses campos sao a variavel independente<br/>do experimento. No cliente, o estudante saberia<br/>de que lado o algoritmo o colocou.
  S->>P: COMMIT (solta o lock)
  S-->>C: posts, modulos, diagnostico
  C-->>E: so posts e modulos (o diagnostico fica no servidor)
\end{lstlisting}

% ---------------------------------------------------------------------------
\section{Figura~\ref{fig:estados-geracao} --- Estados de uma geração assíncrona}

\begin{lstlisting}[language=mermaid]
stateDiagram-v2
  [*] --> pendente : POST (202)
  pendente --> gerando : claim atomico (updateMany)
  gerando --> pronto : gravou o plano
  gerando --> erro : falha do provider, MAX_TOKENS ou orfa ja na 2a tentativa
  gerando --> gerando : orfa (>12 min sem worker) e so 1 tentativa - novo claim, disparado pela leitura do plano
  pronto --> [*]

  note right of pendente
    Plano de aula (app-90, planos/). O claim atomico
    (updateMany: pendente, ou gerando ha mais de 12 min)
    substitui uma fila externa:
    duas instancias podem tentar, so uma ganha a linha.
    Menor peca possivel que resolve o problema (ADR-17).
  end note
\end{lstlisting}

% ---------------------------------------------------------------------------
\section{Figura~\ref{fig:uml-taxonomia} --- Os catorze tipos de diagrama da UML 2.5.1}

\begin{lstlisting}[language=mermaid]
classDiagram
  class Diagrama {
    <<abstract>>
  }
  class Estrutura["Diagrama de estrutura"] {
    <<abstract>>
  }
  class Comportamento["Diagrama de comportamento"] {
    <<abstract>>
  }
  class Interacao["Diagrama de interacao"] {
    <<abstract>>
  }
  Diagrama <|-- Estrutura
  Diagrama <|-- Comportamento

  %% estrutura: sete folhas
  class Classes {
    fig:uml-classes-dominio, fig:uml-classes-ia
    antes: fig:c4l4, fig:graus
  }
  class Objetos {
    fig:uml-objetos
  }
  class Pacotes {
    fig:uml-pacotes
    antes: fig:monorepo (grafo)
  }
  class Componentes {
    fig:uml-componentes
  }
  class EstruturaComposta["Estrutura composta"] {
    fig:uml-estrutura
  }
  class Implantacao["Implantacao"] {
    fig:uml-implantacao
    antes: fig:deploy-contrato (infra)
  }
  class Perfil {
    fig:uml-perfil
  }
  Estrutura <|-- Classes
  Estrutura <|-- Objetos
  Estrutura <|-- Pacotes
  Estrutura <|-- Componentes
  Estrutura <|-- EstruturaComposta
  Estrutura <|-- Implantacao
  Estrutura <|-- Perfil

  %% comportamento: tres folhas e o ramo abstrato da interacao
  class CasosDeUso["Casos de uso"] {
    fig:uml-casos-gestao
    antes: fig:casos-de-uso
  }
  class Atividades {
    fig:uml-atividade-disp
    antes: fig:fluxo-embed, fig:guardrails-bemestar, fig:app92-caminho
  }
  class MaquinaDeEstados["Maquina de estados"] {
    fig:uml-estados-ponte
    antes: fig:estados-geracao
  }
  Comportamento <|-- CasosDeUso
  Comportamento <|-- Atividades
  Comportamento <|-- MaquinaDeEstados
  Comportamento <|-- Interacao

  class Sequencia["Sequencia"] {
    fig:uml-seq-entrada, fig:uml-seq-ia, fig:uml-seq-agente, fig:uml-seq-ajuda
    antes: fig:seq-feed, fig:seq-handshake, fig:seq-erro-cliente
  }
  class Comunicacao["Comunicacao"] {
    fig:uml-comunicacao
  }
  class Tempo {
    fig:uml-tempo
  }
  class VisaoGeral["Visao geral de interacao"] {
    fig:uml-visao-geral
  }
  Interacao <|-- Sequencia
  Interacao <|-- Comunicacao
  Interacao <|-- Tempo
  Interacao <|-- VisaoGeral

  note "Como ler: abstract e classificador abstrato - nao existe diagrama<br/>'de estrutura' que nao seja de um tipo concreto. 'antes' aponta para<br/>a figura do mesmo tipo que o documento ja tinha; (grafo) e (infra)<br/>marcam as que cumprem o papel com outra notacao."
\end{lstlisting}

% ---------------------------------------------------------------------------
\section{Figura~\ref{fig:uml-perfil} --- Perfil UML da plataforma}

\begin{lstlisting}[language=mermaid]
%% aproximacao: o Mermaid nao tem diagrama de perfil; a extensao de metaclasse (seta de ponta cheia) vira associacao rotulada "estende", e as restricoes entre chaves vao para notas.
classDiagram
  %% [x] no corpo da classe = {x} da figura: chave fecha o corpo no Mermaid.
  namespace PerfilUnooria {
    class TabelaDeTenant {
      <<estereotipo>>
      prefixo: String [0..1]
      codigo da experiencia, vazio na plataforma
    }
    class Experiencia {
      <<estereotipo>>
      id: String, o expId
      versao: SemVer
      produto: Produto
      publico: Publico = estudante
      permissoes: String [*]
      viewport: String [0..1]
    }
    class Bundle {
      <<estereotipo>>
      caminho = gs://[contrato]-exp/[expId]/[versao]/
    }
    class RotaDeEstudante {
      <<estereotipo>>
      EXP_ID: String
    }
    class ChamadaDeIA {
      <<estereotipo>>
      perfil: estudante | professor
    }
    class Produto {
      <<enumeracao>>
      entre-redes
      plano-de-aula
      plano-de-ensino
      bem-estar
      agente-professor
      ajuda
    }
    class Publico {
      <<enumeracao>>
      estudante
      professor
    }
  }

  %% metaclasses da UML, fora do perfil
  class MClass["Class"] {
    <<metaclasse>>
  }
  class MComponent["Component"] {
    <<metaclasse>>
  }
  class MArtifact["Artifact"] {
    <<metaclasse>>
  }
  class MOperation["Operation"] {
    <<metaclasse>>
  }

  TabelaDeTenant --> MClass : estende
  Experiencia --> MComponent : estende
  Bundle --> MArtifact : estende
  RotaDeEstudante --> MOperation : estende
  ChamadaDeIA --> MOperation : estende
  Bundle ..> Experiencia : manifesta
  Experiencia ..> Produto

  note for TabelaDeTenant "{escola_id NOT NULL}<br/>{ENABLE e FORCE ROW LEVEL SECURITY}<br/>{toda query dentro de comTenancy}"
  note for Experiencia "{gabarito e correcao so na API}<br/>{token so em memoria}"
  note for Bundle "{imutavel: mudar o conteudo e subir a versao}<br/>{versao igual ao version do package.json}"
  note for RotaDeEstudante "{chama exigirMissaoLiberada no funil}<br/>{403 sem linha em liberacoes}"
  note for ChamadaDeIA "{QuotaService.temSaldo antes}<br/>{QuotaService.consumir depois}<br/>{vendor so atras de uma porta de ai/: AI_PROVIDER ou AGENTE_PROFESSOR}"
  note "Como ler: um elemento marcado TabelaDeTenant em outro diagrama<br/>carrega as tres restricoes sem repeti-las. A seta cheia e extensao<br/>de metaclasse; nao confundir com a heranca, de seta aberta."
  note "Onde o perfil ainda nao e lei: geracao/ e as reacoes de motor/<br/>chamam o provider sem temSaldo - as excecoes registradas na ADR-74<br/>e em ai/AGENTS.md. Uma ChamadaDeIA nova nao herda a excecao."
  %% A Ajuda com avatar (ADR-132) nao e ChamadaDeIA: e a excecao unica e
  %% deliberada a quota de tokens, limitada por minutos de video.
\end{lstlisting}

% ---------------------------------------------------------------------------
\section{Figura~\ref{fig:uml-casos-gestao} --- Casos de uso da administração da rede}

\begin{lstlisting}[language=mermaid]
flowchart LR
  gesc([Gestor de escola]); greg([Gestor regional]); gctr([Gestor do contrato]); gadm(["Admin (Unooria)"])
  prof([Professor]); ci[[CI/CD]]; visit([Visitante do site])

  %% generalizacao de atores (triangulo aberto na figura): o filho faz tudo o que o pai faz
  greg ==>|generalizacao| gesc
  gctr ==>|generalizacao| greg
  gadm ==>|generalizacao| gctr

  subgraph SIS["Plataforma Unooria - administracao"]
    UC20(("UC-20 Importar pessoas e enturmar"))
    UC29(("UC-29 Cadastrar pessoa no balcao"))
    UC21(("UC-21 Editar pessoa e redefinir senha"))
    UC30(("UC-30 Conferir a cadeia de atribuicao"))
    UC32(("UC-32 Trocar a escola ativa"))
    UC33(("UC-33 Ver o panorama da diretoria"))
    UC34(("UC-34 Gerar o relatorio da diretoria"))
    UC31(("UC-31 Criar papel sem escola"))
    UC35(("UC-35 Manter escolas e regionais"))
    UC22(("UC-22 Ler o agregado da rede"))
    UC36(("UC-36 Ler o historico de recursos"))
    UC37(("UC-37 Ligar ou desligar recurso"))
    UC39(("UC-39 Editar a quota de IA"))
    UC40(("UC-40 Ler cadastros de interesse"))
    UC24(("UC-24 Ativar ou reverter versao"))
    UC38(("UC-38 Fixar versao por escopo"))
    UC42(("UC-42 Liberar ou encerrar missao na turma"))
    UC43(("UC-43 Conferir a disponibilidade"))
    UC44(("UC-44 Montar as equipes da missao"))
    UC23(("UC-23 Publicar manifesto"))
    UC41(("UC-41 Cadastrar interesse"))
  end

  gesc --- UC20 & UC29 & UC21
  greg --- UC32 & UC33 & UC34
  gctr --- UC31 & UC35 & UC22 & UC36
  gadm --- UC37 & UC39 & UC40
  prof --- UC42 & UC44
  ci --- UC23
  visit --- UC41

  UC29 -.->|include| UC30
  UC29 -.->|include| UC20
  UC31 -.->|include| UC30
  UC42 -.->|include| UC43
  UC24 -.->|extend| UC37
  UC38 -.->|extend| UC37
  UC24 -.->|extend| UC23

  %% caso de uso de suporte (incluido ou extensao)
  classDef suporte stroke-dasharray: 4 3
  class UC30,UC24,UC38,UC43 suporte
\end{lstlisting}

% ---------------------------------------------------------------------------
\section{Figura~\ref{fig:uml-classes-dominio} --- Classes: o modelo de domínio da plataforma}

\begin{lstlisting}[language=mermaid]
classDiagram
  %% [x] no corpo da classe = {x} da figura: chave fecha o corpo no Mermaid.
  %% As restricoes entre chaves de cada classe estao nas notas.
  class Escola {
    nome: String
    codigoInep: String [0..1]
    tipo: TipoDeEscola
    municipio, uf: String [0..1]
  }
  class Turma {
    nome: String
    serie: String
    tipo: regular | componente
    turno: String [0..1]
  }
  class PeriodoLetivo {
    tipo, titulo: String
    ano: Int
    inicio, fim: Date
  }
  class Lotacao {
    papel: estudante | professor | gestor_escola
    principal: Boolean
    inicio: Date
    fim: Date [0..1]
  }
  class Matricula {
    <<TabelaDeTenant>>
    papel: estudante | professor
    principal: Boolean
    inicio: Date
    fim: Date [0..1]
  }
  class Liberacao {
    <<TabelaDeTenant>>
    estado: liberada | encerrada
    liberadaEm: DateTime
    encerradaEm [0..1]
  }
  class Usuario {
    username: String
    displayName: String
    role: Papel
    ativo: Boolean
    senhaProvisoria: Boolean
    /escola: Escola [0..1]
    /turma: Turma [0..1]
  }
  class Disponibilidade {
    ligado: Boolean
    versaoFixa: String [0..1]
    motivo: String [0..1]
    definidoEm: DateTime
  }
  class Experiencia {
    expId: String
    versaoAtiva: String
    series: String [*]
    publico: Publico
    permissoes: String [*]
  }
  class Papel {
    <<enumeracao>>
    estudante
    professor
    gestor_escola
    gestor_regional
    gestor_contrato
    admin
  }
  class Recurso {
    <<abstract>>
    /recursoId: String
    produto:[slug] ou experiencia:[expId]
  }
  class Produto {
    slug: Produto
    vocabulario em recursos/catalogo.ts
  }

  Escola "1" -- "*" Turma
  Turma "*" -- "1" PeriodoLetivo
  Escola "pai 0..1" -- "filhas *" Escola
  Escola "1" -- "*" Lotacao
  Lotacao "*" -- "1" Usuario
  Turma "1" -- "*" Matricula
  Matricula "*" -- "1" Usuario
  Turma "1" -- "*" Liberacao
  Liberacao "*" -- "1" Experiencia : por expId, sem FK
  Usuario -- "0..1" Escola : regional
  Disponibilidade "*" -- "0..1" Escola : escopo
  Disponibilidade "*" -- "1" Recurso
  Recurso <|-- Produto
  Recurso <|-- Experiencia
  Produto "0..1" -- "*" Experiencia : teto
  Usuario ..> Papel

  note for Escola "{regional e Escola com tipo = regional}<br/>{o contrato nao e linha: e a stack}"
  note for Turma "{unica por escola, periodo e nome}<br/>{fora da RLS: o filtro e a tranca}"
  note for PeriodoLetivo "{do contrato, sem escola}"
  note for Lotacao "{unica por escola, usuario e papel}"
  note for Matricula "{no maximo uma principal ativa por estudante}"
  note for Liberacao "{sem linha = bloqueada}<br/>{unica por turma e expId}"
  note for Usuario "{escola obrigatoria, salvo papeis sem escola}<br/>{role = gestor_regional se e so se regional nao nula}"
  note for Disponibilidade "{sem linha = ligado, salvo produto em teste interno}<br/>{unica por recurso e escopo}<br/>{escopo nulo = o contrato}"
  %% Teste interno = PRODUTOS_DESLIGADOS_POR_PADRAO (recursos/catalogo.ts), hoje
  %% so a ajuda (ADR-132). O padrao mora no resolvedor, nao numa linha semeada.
\end{lstlisting}

% ---------------------------------------------------------------------------
\section{Figura~\ref{fig:uml-classes-ia} --- Classes: a porta da IA}

\begin{lstlisting}[language=mermaid]
classDiagram
  %% [x] no corpo da classe = {x} da figura: chave fecha o corpo no Mermaid.
  %% Linha sem parenteses cai no compartimento de atributos: as anotacoes
  %% cinza que a figura poe sob as operacoes aparecem aqui entre os atributos.
  class AiController {
    <<controller>>
    -provider: AiProvider
    -quota: QuotaService
    +gerar(corpo) [texto]
    POST /ai/gerar, so OidcGuard
  }
  class PlanosService {
    <<service>>
    e mais treze modulos de dominio
    +gerar(...)
    temSaldo -> provider -> consumir
  }
  class AgenteProfessorService {
    <<service>>
    -agente: AgenteProfessorPort
    -quota: QuotaService
    +temSaldo(quem)
    +transmitir(...) SSE
  }
  class AiResult {
    <<tipo de dado>>
    texto: string
    tokensUsados: number
    tokensEntrada?, tokensSaida?
    finishReason?
    voltas?, tentativas?
  }
  class QuotaService {
    -prisma: PrismaService
    +limite(perfil) Promise~number~
    +usadoHoje(escolaId, usuarioId) Promise~number~
    +temSaldo(escolaId, usuarioId, perfil = "estudante") Promise~boolean~
    +consumir(escolaId, usuarioId, tokens) Promise~number~
  }
  class AiProvider {
    <<interface>>
    injetado pelo token AI_PROVIDER
    +generate(req) Promise~AiResult~
    +generateJson~T~(req) Promise~T~
    +generateImage(prompt, opcoes?) Promise~Buffer~
    +generateSpeech(req) Promise~Buffer~
    +conversar(req) Promise~AiResult~
    +conversarEmFluxo(req) AsyncIterable
  }
  class AgenteProfessorPort {
    <<interface>>
    +modo: ModoDoAgente
    injetada pelo token AGENTE_PROFESSOR
    +criarSessao(...)
    +lerSessao(...)
    +apagarSessao(...)
    +renovarSessao?(...)
    +enviar(...) AsyncIterable
  }
  class AiModule {
    <<modulo @Global>>
    exports: [AI_PROVIDER, QuotaService, RegistroDeConsumo, AGENTE_PROFESSOR, AVATAR_AO_VIVO, VOZ_AO_VIVO]
    +useFactory(config, metricas, registro) AiProvider
    +useFactory(config) AgenteProfessorPort
    = criarAgenteProfessor(modo, ...)
  }
  class VertexProvider {
    -auth: GoogleAuth
    ADC, sem chave em env
    as seis operacoes, por :generateContent; o fluxo, por :streamGenerateContent
  }
  class MockProvider {
    deterministico
    generate, conversar, conversarEmFluxo: respondem
    os outros tres: lancam
  }
  class MedidorDeProvider {
    <<decorador>>
    -real: AiProvider
    -metricas: Metricas
    -modelos: ModelosDoMedidor
    as seis operacoes, cada uma por medir(op, modelo, fn)
  }
  class ClienteAgentEngine {
    modo agent-engine (ADC) ou local (shim)
    :query e :streamQuery?alt=sse no Agent Engine
  }
  class ErroDeGeracao {
    extends Error
  }
  class ia_quotas {
    perfil PK
    tokens_por_dia
    sem escola: do contrato
  }
  class ia_uso_diario {
    <<TabelaDeTenant>>
    escola_id, usuario_id, dia PK
    tokens_usados
  }

  VertexProvider ..|> AiProvider
  MockProvider ..|> AiProvider
  MedidorDeProvider ..|> AiProvider
  MedidorDeProvider o-- "1" AiProvider : real
  ClienteAgentEngine ..|> AgenteProfessorPort
  MockProvider ..> ErroDeGeracao : lanca
  AiProvider ..> AiResult : usa
  AiController ..> QuotaService
  AiController ..> AiProvider
  PlanosService ..> AiProvider : @Inject(AI_PROVIDER)
  PlanosService ..> QuotaService
  AgenteProfessorService ..> AgenteProfessorPort : @Inject(AGENTE_PROFESSOR)
  AgenteProfessorService ..> QuotaService
  AiModule ..> MedidorDeProvider : cria
  AiModule ..> ClienteAgentEngine : cria
  QuotaService ..> ia_quotas : le
  QuotaService ..> ia_uso_diario : upsert

  note for AiModule "A fabrica, no codigo: new MedidorDeProvider(env === 'mock'<br/>? new MockProvider() : new VertexProvider(), metricas, modelos, ...).<br/>Trocar de fornecedor e um arquivo novo em src/ai e uma linha aqui."
  note "O que o tipo nao garante: nada obriga a chamar temSaldo, e entre ele<br/>e o consumir nao ha trava: duas chamadas simultaneas podem passar<br/>o teto em uma resposta cada."
  %% AgenteProfessorPort tambem e implementada por AgenteProfessorMock (suite e2e)
  %% e AgenteProfessorDesligado (o padrao), no mesmo agent-engine.client.ts.
\end{lstlisting}

% ---------------------------------------------------------------------------
\section{Figura~\ref{fig:uml-objetos} --- Objetos: um instantâneo da disponibilidade}

\begin{lstlisting}[language=mermaid]
%% aproximacao: o Mermaid nao tem diagrama de objetos; cada instancia vira um no com os valores dos slots, o nome sublinhado vira a tag u do HTML, e cada link vira aresta rotulada com o papel.
flowchart TB
  norte["<u>norte : Escola</u><br/>tipo = 'regional'<br/>nome = 'Regional Norte'"]
  ef["<u>ef : Escola</u><br/>tipo = 'escola'<br/>nome = 'EE Cora Coralina'"]
  ee["<u>ee : Escola</u><br/>tipo = 'escola'<br/>nome = 'EE Monteiro Lobato'"]

  d3["<u>d3 : Disponibilidade</u><br/>recursoId =<br/>'experiencia:exp06ef2b'<br/>ligado = false<br/>motivo = 'adiar'"]
  d4["<u>d4 : Disponibilidade</u><br/>recursoId =<br/>'produto:plano-de-aula'<br/>ligado = false<br/>motivo = 'nao contratado'"]
  d1["<u>d1 : Disponibilidade</u><br/>recursoId =<br/>'experiencia:exp06ef2b'<br/>escola = nulo (contrato)<br/>ligado = false"]
  d2["<u>d2 : Disponibilidade</u><br/>recursoId =<br/>'experiencia:exp06ef2b'<br/>ligado = true<br/>motivo = 'piloto'"]

  t6a["<u>t6a : Turma</u><br/>nome = '6A'<br/>serie = '6EF'"]
  paula["<u>aula : Produto</u><br/>slug = 'plano-de-aula'"]
  e1["<u>e1 : Experiencia</u><br/>expId = 'exp06ef2b'<br/>produto = 'entre-redes'<br/>series = ['6EF']"]
  l1["<u>l1 : Liberacao</u><br/>estado = 'liberada'"]
  e2["<u>e2 : Experiencia</u><br/>expId =<br/>'app-90-plano-de-aula'<br/>publico = 'professor'"]

  ef ---|pai| norte
  ee ---|pai| norte
  d3 ---|escopo| ef
  d4 ---|escopo| ee
  d2 ---|escopo| norte
  t6a ---|escola| ee
  l1 ---|turma| t6a
  l1 ---|expId| e1
  d1 ---|recurso| e1
  d2 ---|recurso| e1
  d3 ---|recurso| e1
  d4 ---|recurso| paula
  e2 ---|"produto (teto)"| paula

  %% d1 nao tem link de escopo: escola nula e o contrato inteiro.
\end{lstlisting}

% ---------------------------------------------------------------------------
\section{Figura~\ref{fig:uml-pacotes} --- Pacotes: quem pode importar quem}

\begin{lstlisting}[language=mermaid]
%% aproximacao: o Mermaid nao tem diagrama de pacotes; cada pacote vira subgraph, e a dependencia de codigo (import/use) vira seta tracejada rotulada.
flowchart TB
  subgraph MONO["monorepo - plataforma-unooria"]
    subgraph API["apps/api - @edu/api - NestJS + Prisma"]
      subgraph LTR["transversal"]
        ltr["auth/ tenancy/ observabilidade/ config/"]
      end
      subgraph LPL["plataforma"]
        lpl["usuarios/ rostering/ salas/ equipes/ evidencias/<br/>recursos/ experiencias/ telemetry/ analise/ comum/ ..."]
      end
      subgraph LIN["infraestrutura"]
        lin["prisma/ ai/ object-storage/ assets/"]
      end
      subgraph LDO["dominio {nenhum modulo importa outro}"]
        ldo["motor/ geracao/ vila/ unoogpt/ quem-mandou-isso/ ..."]
      end
    end
    subgraph HUB["apps/hub"]
      hub["@edu/hub<br/>Next.js, App Router<br/>rotas por papel, BFF, JWKS,<br/>montagem do embed"]
    end
    subgraph AGP["apps/agente-professor"]
      agp["Python + Google ADK, no Vertex AI Agent Engine<br/>nao importa nem e importado:<br/>so a API o chama, por HTTP"]
    end
    subgraph SDK["packages/sdk"]
      sdk["@edu/sdk<br/>TypeScript + Zod<br/>ponte postMessage v2, cliente<br/>HTTP, telemetria em lote"]
    end
    subgraph UI["packages/ui"]
      ui["@edu/ui<br/>as oito pontes do DS, uma<br/>por secao: 'use client',<br/>nomes sem colisao"]
    end
    subgraph EXPS["apps/experiencias"]
      exps["_template exp01ef1b ... exp08ef4b<br/>app-70 app-90 app-91 app-92<br/>13 bundles Vite publicados;<br/>_template e backup-* ficam fora do deploy"]
    end
  end
  subgraph DS["externo - npm"]
    ds["@unooria/design-system<br/>repo storybook-unooria<br/>componentes + CSS de secao"]
  end

  %% dentro da api: toda faixa usa a transversal
  ldo -.->|import| lpl
  ldo -.->|import| lin
  lpl -.->|import| lin
  lpl -.->|"import - toda faixa usa a transversal"| ltr

  %% entre pacotes do workspace
  hub -.->|import| sdk
  hub -.->|import| ui
  hub -.->|use: so CSS da secao e fontes| ds
  exps -.->|import| sdk
  exps -.->|import| ui
  exps -.->|use: so CSS da secao e fontes| ds
  ui -.->|import| ds

  %% A api NAO importa nada do workspace (nem @edu/sdk): o contrato da ponte e
  %% duplicado de proposito, com Zod de cada lado. Hub e api nao se importam:
  %% falam por HTTP, pelo BFF. O agente-professor (Python) tambem fica fora do
  %% grafo de import: so a api o chama, por HTTP. A app-70 importa so a SDK
  %% (jogo em canvas, sem DOM de interface). No dominio, duas experiencias
  %% conversam por tabela de contrato ou por endpoint de fronteira, nunca por import.
\end{lstlisting}

% ---------------------------------------------------------------------------
\section{Figura~\ref{fig:uml-componentes} --- Componentes e interfaces}

\begin{lstlisting}[language=mermaid]
%% aproximacao: o Mermaid nao tem diagrama de componentes; cada interface vira um no circular, a linha cheia sai de quem a FORNECE (a bola) e a seta tracejada sai de quem a REQUER (o soquete).
flowchart TB
  hub["componente<br/>Hub<br/>[Next.js]<br/>sessao, papel, embed,<br/>assinatura do token"]
  subgraph EXP["componente - Experiencia [Vite + React, bundle]"]
    sdk["biblioteca<br/>@edu/sdk - ponte, HTTP, telemetria"]
  end
  api["componente<br/>API<br/>[NestJS + Prisma]<br/>autorizacao por papel, tenancy, regra de dominio, gabarito, porta da IA"]
  pg["externo<br/>PostgreSQL<br/>[Cloud SQL]"]
  ps["externo<br/>Pub/Sub"]
  vx["externo<br/>Vertex AI"]
  gcs["externo<br/>Cloud Storage"]
  lav["externo<br/>LiveAvatar<br/>[HeyGen]"]
  agp["externo<br/>Agente do professor<br/>[ADK no Agent Engine]"]

  %% hub e experiencia: a ponte (o hub fornece, a experiencia requer)
  iPonte(("Ponte postMessage v2<br/>hub:init, hub:token,<br/>hub:pause/resume"))
  hub --- iPonte
  EXP -.-> iPonte

  %% hub e api
  iLogin(("POST /auth/login<br/>Credentials"))
  iBff(("REST pelo BFF<br/>/api/plataforma/*"))
  iJwks(("JWKS<br/>/.well-known/jwks.json"))
  api --- iLogin
  hub -.-> iLogin
  api --- iBff
  hub -.-> iBff
  hub --- iJwks
  api -.-> iJwks

  %% experiencia (pela sdk) e api
  iRest(("REST com Bearer<br/>rotas de dominio"))
  iTel(("Telemetria<br/>/telemetry/batch<br/>/telemetry/beacon"))
  api --- iRest
  sdk -.-> iRest
  api --- iTel
  sdk -.-> iTel

  %% api e externos
  iSql(("SQL + RLS<br/>comTenancy"))
  iPub(("Publicacao<br/>{contrato}-telemetria"))
  iIngest(("Ingestao push<br/>/telemetry/ingest"))
  iGen(("Geracao e voz ao vivo<br/>ADC, :generateContent,<br/>Live por WebSocket"))
  iAgp(("Turnos do agente<br/>:query, :streamQuery"))
  iObj(("Objetos<br/>assets, planos"))
  pg --- iSql
  api -.-> iSql
  ps --- iPub
  api -.-> iPub
  api --- iIngest
  ps -.-> iIngest
  vx --- iGen
  api -.-> iGen
  gcs --- iObj
  api -.-> iObj
  agp --- iAgp
  api -.-> iAgp
  %% So a API chama o Agente do professor; o navegador nunca fala com o Agent Engine.

  %% api e LiveAvatar (ADR-132): dois sentidos
  iLlm(("LLM<br/>/ajuda/llm"))
  iSess(("Sessoes<br/>/v1/sessions"))
  api --- iLlm
  lav -.-> iLlm
  lav --- iSess
  api -.-> iSess
  %% O video vai do vendor direto ao navegador, sem interface com a API.

  %% O Cloud Storage tambem serve o bundle ao navegador pela CDN: e artefato,
  %% nao interface, e esta no diagrama de implantacao (fig:uml-implantacao).
\end{lstlisting}

% ---------------------------------------------------------------------------
\section{Figura~\ref{fig:uml-estrutura} --- Estrutura composta: a página da missão}

\begin{lstlisting}[language=mermaid]
%% aproximacao: o Mermaid nao tem diagrama de estrutura composta; as partes viram subgraphs aninhados e nos, as portas viram nos pequenos dentro da parte a que pertencem, e a fronteira do sandbox e o subgraph do iframe.
flowchart LR
  page["server component<br/>app/exp/[expId]/page.tsx<br/>montarEmbed(): token,<br/>src, expOrigin,<br/>permissoes, viewport"]
  cdn["CDN<br/>{cdn}/{expId}/{versao}/<br/>index.html?hubOrigin=..."]
  rprog["BFF<br/>POST /api/plataforma/<br/>progresso/:expId/<br/>concluir"]
  rtok["route handler<br/>GET /api/exp-token<br/>401/403 -> hub:pause<br/>e o aviso"]
  rlog["route handler<br/>POST /api/ponte<br/>uma linha de log; 204"]

  subgraph EXPHOST["ExpHost - client component"]
    pEmbed["porta"]
    pProg["porta"]
    subgraph DISP["dispositivo : DispositivoFrame [0..1] - so quando o manifesto declara viewport"]
      subgraph FRAME["frame : ExpFrame [1]"]
        pTok["porta"]
        pLog["porta"]
        proto["ponte : protocolo v2<br/>postMessage;<br/>ev.origin + Zod"]
        aviso["aviso : AlertBox [0..1]<br/>'Sua sessao expirou'; substitui o iframe"]
        subgraph IFRAME["iframe [1] - fronteira do sandbox"]
          pSrc["porta src"]
          pPonte["porta ponte"]
          sbx["um HTMLIFrameElement com<br/>sandbox = 'allow-scripts<br/>allow-same-origin allow-modals'<br/>+ pointer-lock, forms, downloads<br/>se o manifesto pedir<br/>allow = 'autoplay; fullscreen;<br/>clipboard-write'"]
          bexp["exp : Experiencia (estereotipo)<br/>o bundle, servido pela CDN"]
          bsdk["sdk : @edu/sdk"]
          btok["token : string<br/>{so em memoria}"]
        end
      end
    end
  end

  %% conectores; === marca os dois unicos que atravessam o sandbox
  page ---|embed| pEmbed
  pProg ---|progresso| rprog
  pTok ---|token| rtok
  pLog ---|sinal| rlog
  pPonte === proto
  proto --- pTok
  proto --- pLog
  cdn === pSrc
  pPonte --- bexp
  bexp --- bsdk
  bsdk --- btok

  classDef borda stroke:#e8740c,stroke-width:2px
  class pPonte,pSrc borda

  %% targetOrigin e sempre expOrigin, nunca *: o frame so envia para a origem
  %% que montarEmbed calculou, e so aceita mensagem cujo ev.origin seja ela.
  %% O ExpHost tem uma segunda sobreposicao, fora do desenho: 'A missao
  %% travou', com o hash do erro, quando a experiencia manda exp:erro. E dele,
  %% nao do frame: sobrevive a remontagem do frame ao recarregar.
\end{lstlisting}

% ---------------------------------------------------------------------------
\section{Figura~\ref{fig:uml-atividade-disp} --- Atividades: o resolvedor de disponibilidade}

\begin{lstlisting}[language=mermaid]
flowchart TB
  subgraph ESQ["comum/disponibilidade.ts - com banco"]
    ini([inicio])
    a1["buscar experiencias[expId]<br/>produto, publico; pode nao existir"]
    a2["cadeiaDaEscola<br/>[escola] ou [escola, regional]"]
    a3["linhasDoEscopo<br/>escola nula OU na cadeia; so ids legiveis"]
    a4["produtoConhecido<br/>fora do vocabulario -> null"]
    d4{"efetivo.ligado?"}
    fim(((fim)))
    a7["lanca 403 recurso-indisponivel"]
    ff((("X - fim de fluxo")))
  end

  subgraph DIR["recursos/regras.ts - funcao pura"]
    d1{"produto?"}
    a5["maisEspecifica(<br/>'produto:{slug}')<br/>sem linha: ligadoPorPadrao"]
    d2{"produto ligado?"}
    r1["desligado,<br/>tetoDoProduto"]
    m1{" "}
    a6["maisEspecifica(<br/>'experiencia:{expId}')"]
    d3{"linha da experiencia?"}
    r2["ligado,<br/>versaoFixa = null"]
    r3["ligado e versaoFixa<br/>da linha"]
    m2{" "}
  end

  ini --> a1 --> a2 --> a3 --> a4
  a4 -->|resolverExperiencia| d1
  d1 -->|sim| a5
  a5 --> d2
  d2 -->|nao| r1
  d2 -->|sim| m1
  d1 -->|nao| m1
  m1 --> a6 --> d3
  d3 -->|nenhuma| r2
  d3 -->|uma linha| r3
  r2 --> m2
  r3 --> m2
  r1 --> m2
  m2 --> d4
  d4 -->|sim| fim
  d4 -->|nao| a7
  a7 --> ff

  %% m1 e m2 sao nos de juncao. O X e o fim de fluxo pela excecao: a transacao
  %% de quem chamou desfaz o que tiver feito.
  %% maisEspecifica percorre a cadeia na ordem - escola, depois regional - e
  %% devolve a primeira linha do recurso; sem nenhuma, a de escola nula (o
  %% contrato); sem ela, null. Sobe so um nivel: a rede nao tem arvore mais funda.
  %% A linha inteira e a decisao: ligado e versaoFixa nao se resolvem em cadeias
  %% separadas. A tela mostra exatamente a linha que vale.
  %% Sem linha em experiencias (stack nova, antes do primeiro publish) NAO tranca:
  %% sem produto, vale so o interruptor da propria experiencia.
  %% A mensagem depende do publico: 'Este recurso nao esta disponivel para a sua
  %% rede.' para professor; 'Esta missao nao esta aberta agora.' para estudante.
  %% Produto sem linha vale o padrao dele (ligadoPorPadrao): ligado, salvo os de
  %% PRODUTOS_DESLIGADOS_POR_PADRAO (hoje so a ajuda, ADR-132). A experiencia
  %% sem linha continua ligada.
\end{lstlisting}

% ---------------------------------------------------------------------------
\section{Figura~\ref{fig:uml-seq-entrada} --- Sequência: a entrada do estudante na missão}

\begin{lstlisting}[language=mermaid]
sequenceDiagram
  participant N as Navegador (hub no cliente)
  participant H as Hub (servidor)
  participant A as API (NestJS)
  participant P as PostgreSQL (Cloud SQL)

  N->>H: GET /exp/:expId (cookie)
  H->>A: GET /salas/minhas
  A->>P: turmasDoEstudante + liberacoes
  P-->>A: matriculas e linhas
  A-->>H: grade: estado por missao, sem as indisponiveis

  alt missao nao esta liberada na grade
    H-->>N: 404 (notFound)
  else liberada
    H->>H: montarEmbed: assina o token (RS256, 15 min)
    H->>A: GET /experiencias/:expId e /recursos/efetivos
    A-->>H: manifesto, versao fixada da escola
    H-->>N: HTML com ExpFrame
  end

  %% a linha de vida da experiencia comeca na mensagem create: antes dela, o iframe nao existe
  create participant E as Experiencia (iframe)
  N-->>E: create - bundle da CDN

  Note over N,E: ref: handshake exp:ready / hub:init - fig:seq-handshake

  E->>A: 1a rota de estudante (Bearer)
  A->>A: funil: garantirSessao ou estudanteDaEscola
  A->>P: exigirRecursoDisponivel (fig:uml-atividade-disp)
  A->>P: missoesLiberadasPara
  P-->>A: linhas de disponibilidade e de liberacao

  alt recurso desligado no escopo que vale
    A-->>E: 403 recurso-indisponivel
  else sem liberacao liberada nas turmas do estudante
    A-->>E: 403 missao-nao-liberada
  else as duas passam
    A-->>E: 200 - a rota segue
  end

  %% O primeiro alt (o filtro do hub) e UX; as travas estao no segundo, na API,
  %% e disparam em serie: disponibilidade antes de liberacao.
\end{lstlisting}

% ---------------------------------------------------------------------------
\section{Figura~\ref{fig:uml-seq-ia} --- Sequência: uma chamada de IA com quota}

\begin{lstlisting}[language=mermaid]
sequenceDiagram
  participant E as Experiencia (pela SDK)
  participant G as OidcGuard (auth/)
  participant C as AiController (ai/)
  participant Q as QuotaService (ai/)
  participant P as PostgreSQL (Cloud SQL)
  participant M as provider (MedidorDeProvider)
  participant V as Vertex AI (Gemini)

  E->>G: POST /ai/gerar (Bearer)
  G->>G: TokenService.verificar: JWKS do hub
  break token ausente ou invalido
    G-->>E: 401
  end
  G->>C: req.auth = {usuarioId, escolaId, perfil}
  C->>C: GerarSchema.parse (Zod, 400 se invalido)
  C->>Q: temSaldo(escola, usuario, perfil)
  Q->>P: limite: ia_quotas, fora do comTenancy
  P-->>Q: tokens_por_dia
  Q->>P: usadoHoje: ia_uso_diario, em comTenancy
  P-->>Q: tokens_usados
  Q-->>C: usado < limite
  break sem saldo
    C-->>E: 429 'Quota diaria de IA atingida'
  end
  C->>C: promptDoTutor(expId)
  C->>P: serieDoPedido: turma.findFirst
  P-->>C: serie -> '6o ano'
  C->>M: generate({system, user})
  M->>M: medir: VertexProvider
  M->>V: :generateContent (ADC)
  V-->>M: texto e uso
  M-->>C: AiResult {texto, tokensUsados}
  C->>Q: consumir(..., tokensUsados)
  Q->>P: upsert em comTenancy
  P-->>Q: tokens_usados
  C-->>E: 200 {texto}

  %% break e a saida antecipada: quando a guarda vale, o operando roda e o resto
  %% da interacao NAO acontece. O consumo e gravado DEPOIS da resposta do modelo,
  %% com os tokens que ele declarou; entre o temSaldo e o consumir nao ha trava.
\end{lstlisting}

% ---------------------------------------------------------------------------
\section{Figura~\ref{fig:uml-seq-agente} --- Sequência: um turno do agente de IA do professor}

\begin{lstlisting}[language=mermaid]
sequenceDiagram
  participant N as Navegador (painel do chat)
  participant H as Hub (rota de stream)
  participant A as AgenteProfessor (controller + service)
  participant P as PostgreSQL (Cloud SQL)
  participant E as Agente (ADK no Agent Engine)
  participant G as Gemini (global -> us)
  participant D as Data stores (Vertex AI Search)

  N->>H: POST {texto}
  H->>A: POST .../mensagens (token curto)
  A->>P: daPosse (dono e escola)
  P-->>A: sessao_id
  A->>P: temSaldo (perfil professor)
  P-->>A: usado < limite
  break alguma recusa
    A-->>N: 4xx em JSON, antes do stream
  end
  A-->>N: 200 event-stream (repassado pelo hub)
  A->>E: :streamQuery?alt=sse user_id, session_id
  E->>E: guarda de escopo (modelo leve)
  E->>G: gerar, em stream
  G-->>E: function_call
  E->>D: buscar_habilidades ou buscar_acervo
  D-->>E: habilidades ou trechos
  E->>G: com o retorno
  G-->>E: texto em pedacos
  loop cada evento do agente
    E-->>A: evento: parcial, final ou state_delta
    A-->>N: event: delta (sem repetir o final agregado)
  end
  alt o agente terminou
    A-->>N: event: fontes
    A->>P: atualizada_em, titulo
    A-->>N: titulo, aviso, fim
  else falha do agente ou 240 s
    A-->>N: event: erro
  end
  A->>P: finally: consumir (soma dos nao parciais)
  A->>E: PATCH expireTime (so depois de fim)

  %% As recusas do break, todas em JSON: 401 sem token, 403 sem o papel
  %% professor ou com o produto desligado (OidcGuard, RolesGuard, @Produto),
  %% 400 corpo invalido, 404 conversa de outro dono, 429 RATE_LIMIT (10/min) ou
  %% QUOTA. O erro do stream e AGENTE_INDISPONIVEL ou TEMPO_ESGOTADO (240 s).
  %% Durante o loop a api escreve ": ping" a cada 15 s. "Parar" (ou fechar a
  %% aba) aborta cada salto ate o Vertex; fechar o PAINEL nao aborta. O
  %% consumo roda no finally tambem no abandono e no erro; a renovacao so depois
  %% de fim, e a falha dela e so log.
\end{lstlisting}

% ---------------------------------------------------------------------------
\section{Figura~\ref{fig:uml-seq-ajuda} --- Sequência: uma sessão da Ajuda com avatar}

\begin{lstlisting}[language=mermaid]
sequenceDiagram
  participant N as Navegador (modal da Ajuda)
  participant A as ajuda/ (controllers + service)
  participant P as PostgreSQL (Cloud SQL)
  participant V as LiveAvatar (HeyGen, FULL Mode)
  participant G as Gemini (pelo AiProvider)

  N->>A: POST /ajuda/sessoes (BFF)
  A->>P: fecha as abertas; le ajuda_limites e ajuda_sessoes
  P-->>A: segundos usados no dia (fuso da plataforma)
  break desligado, sem atendente, mais de 6/min ou menos de 1 min no dia
    A-->>N: 403 | 503 | 429, antes da sala
  end
  A->>P: reserva (inicio, duracao_max_s)
  A->>A: assina o token da sessao (HS256)
  A->>V: POST /v1/sessions/token (FULL, com o token da sessao)
  V-->>A: session_id, session_token
  A-->>N: 201 {id, tokenDaSessao, duracaoMaximaS, ...}
  N->>V: start(): entra na sala (WebRTC)
  V->>A: GET /ajuda/llm/models (Bearer)
  loop cada fala, por voz ou por escrito
    N->>V: fala ou message(texto)
    V->>A: POST /ajuda/llm/chat/completions (Bearer + [sessao:...] no system)
    A->>A: duas trancas (senao 401); produto ligado? ate 20/min? (senao vira fala)
    A->>G: conversarEmFluxo (prefixo fixo + papel, historico)
    G-->>A: pedacos (SSE do Vertex)
    A-->>V: chat.completion.chunk em SSE, ate [DONE]
    V-->>N: voz, video e transcricoes, pela sala
  end
  N->>V: stop() (fechar o modal)
  N->>A: POST .../encerrar (keepalive)
  A->>P: fim = agora
  A->>V: POST /v1/sessions/stop

  %% O navegador fala com a API pelo BFF do hub (/api/plataforma/ajuda/**).
  %% Antes do break rodam as guardas da classe: OidcGuard, RolesGuard (nao
  %% estudante) e @Produto("ajuda"). O token da sessao carrega
  %% {sub, escola, perfil, sid} e vai em dynamic_variables.sessao.
  %% Se o vendor nao abre a sala, a linha reservada fecha com duracao zero e a
  %% resposta e 502 AVATAR_FALHOU. No loop, so a autenticacao responde erro: o
  %% vendor fica MUDO diante de status de erro, entao produto desligado e excesso
  %% viram fala em 200, e a falha do modelo vira a fala de desculpa. Sem pedaco
  %% em 7 s sai a fala de espera (o vendor desiste aos 10 s sem o primeiro byte).
  %% O system que o vendor manda e descartado depois de lido o token. O encerrar
  %% escapa do interruptor (@SemDisponibilidade); sessao sem fim conta o teto
  %% inteiro, e abrir outra fecha a anterior. Sem QuotaService: minutos (ADR-132).
\end{lstlisting}

% ---------------------------------------------------------------------------
\section{Figura~\ref{fig:uml-comunicacao} --- Comunicação: o caminho de um lote de telemetria}

\begin{lstlisting}[language=mermaid]
%% aproximacao: o Mermaid nao tem diagrama de comunicacao; cada link vira aresta, a mensagem numerada vira o rotulo dela e a seta da a direcao (o detalhe de cada numero esta na tabela que segue a figura).
flowchart LR
  exp["<u>:Experiencia</u><br/>codigo da missao"]
  tel["<u>t : TelemetryClient</u><br/>@edu/sdk"]
  tok["<u>:TokenService</u><br/>auth/"]
  ctrl["<u>:TelemetryController</u><br/>telemetry/"]
  ps["<u>:Topic</u><br/>{contrato}-telemetria"]
  guard["<u>:PubSubPushGuard</u><br/>ID token do Google"]
  pg["<u>:telemetria_eventos</u><br/>TabelaDeTenant, por mes"]

  exp -->|"1: track()"| tel
  tel -->|"1.1: POST /telemetry/batch<br/>3: sendBeacon /beacon"| ctrl
  ctrl -->|"1.1.1: identidade do token"| ctrl
  ctrl -->|"3.1: verificar()"| tok
  ctrl -->|"1.1.2, 3.2: publishMessage()"| ps
  ps -->|"2: push /ingest"| guard
  guard -->|"2.1: ingest()"| ctrl
  ctrl -->|"2.1.1: INSERT ... DO NOTHING"| pg

  %% Tres gatilhos, tres numeros de topo: 1 e o relogio da SDK, 3 e a saida da
  %% pagina e 2 e o Pub/Sub, quando ele quiser. Nenhum espera o outro.
  %% Sem TELEMETRIA_TOPIC (compose local, CI), 1.1.2 vira uma linha de debug e
  %% nada e gravado.
\end{lstlisting}

% ---------------------------------------------------------------------------
\section{Figura~\ref{fig:uml-tempo} --- Tempo: \emph{token}, renovação e sessão}

\begin{lstlisting}[language=mermaid]
%% aproximacao: o Mermaid nao tem diagrama de tempo; cada faixa vira uma secao do gantt e cada estado um intervalo no eixo de minutos, e as mensagens entre faixas ficam em comentario.
%%{init: {"gantt": {"displayMode": "compact"}}}%%
gantt
  title Missao aberta - o pedido de token dos 20 min falha por rede
  dateFormat mm
  axisFormat %M
  tickInterval 5minute

  section Chamadas da experiencia a API
  200                              :c1, 00, 25m
  401 {5 min de 401}               :crit, c2, 25, 5m
  200                              :c3, 30, 10m

  section Token vigente na experiencia
  A - vale ate 15                  :t1, 00, 10m
  B - vale ate 25 {15 min}         :t2, 10, 15m
  B vencido                        :crit, t2v, 25, 5m
  D - vale ate 45                  :t3, 30, 10m

  section Pedido de token do hub
  pedindo - ok {10 min}            :p1, 10, 36s
  pedindo - falhou (rede)          :crit, p2, 20, 36s
  pedindo - ok                     :p3, 30, 36s
  pedindo                          :p4, 40, 24s

  section Sessao do hub
  ativa {expira 2 h depois do ultimo ponto}  :s1, 00, 40m
  renova o prazo                   :milestone, r0, 00, 0m
  renova o prazo                   :milestone, r1, 10, 0m
  renova o prazo                   :milestone, r2, 30, 0m

  %% Fora dos pulsos, o pedido de token esta ocioso.
  %% Mensagens entre faixas: exp:token-expirando aos 10, 20 e 30 (experiencia
  %% para o hub); hub:token de volta aos 10,6 e aos 30,6 (so os pedidos ok).
  %% O token B (emitido aos 10) vence aos 25; o proximo pedido e aos 30; entre
  %% os dois, toda chamada da experiencia a API recebe 401.
\end{lstlisting}

% ---------------------------------------------------------------------------
\section{Figura~\ref{fig:uml-estados-ponte} --- Máquina de estados de protocolo: a ponte}

\begin{lstlisting}[language=mermaid]
stateDiagram-v2
  [*] --> AguardandoInit : connectToHub(hubOrigin)

  AguardandoInit --> Conectada : hub:init [trunc(v) = 2] / limpa os temporizadores, resolve
  AguardandoInit --> Incompativel : hub:init [trunc(v) != 2]
  AguardandoInit --> TempoEsgotado : after(8 s)

  state Conectada {
    [*] --> Ativa
    Ativa --> Pausada : hub:pause
    Pausada --> Ativa : hub:resume
  }

  Incompativel : Incompativel
  Incompativel : reject - 'Protocolo incompativel'
  TempoEsgotado : TempoEsgotado
  TempoEsgotado : reject - 'Timeout no handshake'
  Incompativel --> [*]
  TempoEsgotado --> [*]

  %% transicoes internas (as que nao mudam de estado): evento [guarda] / efeito
  note right of AguardandoInit
    entry / normalizarHubOrigin, envia exp:ready
    do / reenvia exp:ready a cada 400 ms
    message [origin != hubOrigin] / descarta
    message [nao casa HubInitSchema] / ignora
  end note
  note right of Conectada
    hub:token / token := msg.token
    after(10 min) / envia exp:token-expirando
    hub:init / ignora (resposta a um exp:ready reenviado)
    message [nao casa HubParaExpSchema] / avisa e descarta
  end note

  %% A SDK nao guarda Ativa nem Pausada: repassa as duas mensagens aos handlers
  %% onPause/onResume, e a experiencia para os proprios cronometros. O hub manda
  %% hub:pause em visibilitychange, em blur (salvo quando o foco foi para o
  %% proprio iframe) e antes do aviso de sessao expirada.
  %% Na pratica, o incompativel termina em tempo esgotado: o ExpFrame confere a
  %% versao do exp:ready e, se o maior diverge, registra protocolo-incompativel e
  %% NAO responde. A SDK nunca recebe o hub:init que a levaria a Incompativel:
  %% cai no after(8 s). O estado existe para o dia em que um hub responder com
  %% outra versao.
\end{lstlisting}

% ---------------------------------------------------------------------------
\section{Figura~\ref{fig:uml-visao-geral} --- Visão geral de interação: a jornada do estudante}

\begin{lstlisting}[language=mermaid]
%% aproximacao: o Mermaid nao tem diagrama de visao geral de interacao; cada quadro ref vira um no de sub-rotina com a figura em que a interacao esta desenhada, e as barras de bifurcacao e juncao viram nos estreitos.
flowchart TB
  ini([inicio])
  login[["ref: login (UC-01)<br/>signIn('credentials') -> POST /auth/login"]]
  dsenha{"senha provisoria?"}
  senha[["ref: trocar a senha<br/>/trocar-senha (UC-25)"]]
  m1{" "}
  grade[["ref: a grade do estudante<br/>GET /salas/minhas, sem as indisponiveis"]]
  dlib{"o card esta liberado?"}
  cad["cadeado: 'Ainda nao abriu' ou 'Esta voce ja fechou'"]
  ff1((("X - fim de fluxo")))
  entra[["ref: entrada na missao<br/>as duas travas (fig:uml-seq-entrada)"]]
  hand[["ref: handshake<br/>exp:ready/hub:init (fig:seq-handshake)"]]
  fork["bifurcacao"]
  dom[["ref: rotas de dominio<br/>e IA com quota (fig:uml-seq-ia)"]]
  tel[["ref: telemetria em lote<br/>(fig:uml-comunicacao)"]]
  ren[["ref: renovacao do token<br/>a cada 10 min (fig:uml-tempo)"]]
  join["juncao"]
  dfim{"como a missao termina?"}
  conc[["ref: concluir<br/>POST .../progresso/:expId/concluir"]]
  exp["aviso 'Sua sessao expirou'; hub:pause"]
  trav["'A missao travou', com o hash do erro"]
  f1(((fim))); f2(((fim))); f3(((fim)))

  ini --> login --> dsenha
  dsenha -->|sim| senha
  senha --> m1
  dsenha -->|nao| m1
  m1 --> grade --> dlib
  dlib -->|nao| cad --> ff1
  dlib -->|sim| entra --> hand --> fork
  fork --> dom & tel & ren
  dom & tel & ren --> join
  join --> dfim
  dfim -->|exp:concluir| conc --> f1
  dfim -->|"/api/exp-token da 401 ou 403"| exp --> f2
  dfim -->|exp:erro| trav --> f3

  classDef barra fill:#1f2a44,stroke:#1f2a44,color:#ffffff
  class fork,join barra

  %% As tres interacoes do meio correm juntas enquanto a missao esta aberta. A
  %% saida pelo cadeado termina em fim de fluxo: o estudante continua na home.
\end{lstlisting}

% ---------------------------------------------------------------------------
\section{Figura~\ref{fig:uml-implantacao} --- Implantação: que artefato roda em que nó}

\begin{lstlisting}[language=mermaid]
%% aproximacao: o Mermaid nao tem diagrama de implantacao; cada no (e ambiente de execucao) vira subgraph, cada artefato vira um no dentro dele, o caminho de comunicacao vira linha e a dependencia implanta vira seta grossa.
flowchart LR
  subgraph LAV["externo - LiveAvatar"]
    ilav["HeyGen, na AWS"]
  end
  subgraph NAV["device - Navegador da escola"]
    subgraph NAVENV["browser - aba do hub"]
      apag["artefato<br/>pagina do hub"]
    end
    subgraph IFR["sandbox - iframe"]
      bnav["Bundle<br/>{expId}/{versao}/<br/>token so em memoria"]
    end
  end
  subgraph GHA["no - GitHub Actions"]
    igha["deploy.yml; WIF, sem chave JSON"]
  end

  subgraph GCP["projeto GCP - uma stack por contrato"]
    subgraph LB["no - Load Balancer global + Cloud CDN"]
      ilb["um IP; host rules:<br/>hub, api, exp, assets"]
    end
    subgraph RUNHUB["Cloud Run - {contrato}-frontend"]
      ahub["imagem hub:{sha}<br/>apps/hub/server.js"]
    end
    subgraph JOB["Cloud Run Job - {contrato}-migrate"]
      ajob["imagem api:{sha}<br/>prisma migrate deploy"]
    end
    subgraph GCS["no - Cloud Storage"]
      igcs["-exp, -assets: publicos;<br/>-planos e mapas: PRIVADOS"]
      bgcs["Bundle<br/>{expId}/{versao}/<br/>index.html: max-age=60"]
    end
    subgraph RUNAPI["Cloud Run - {contrato}-api"]
      aapi["imagem api:{sha}<br/>dist/src/main.js"]
    end
    subgraph SQL["no - Cloud SQL {contrato}-pg"]
      isql["PostgreSQL, HA + PITR; so IP privado;<br/>role app sem BYPASSRLS;<br/>a posse das conversas do agente"]
    end
    subgraph PS["no - Pub/Sub"]
      ips["topic -telemetria;<br/>push -telemetria-ingest"]
    end
    subgraph VX["servico gerenciado - Vertex AI"]
      ivx["Gemini em global (reserva us);<br/>Live em us-central1"]
    end
    subgraph ENG["Agent Engine - us-east1 (nasce no apply 2)"]
      aeng["imagem agente-professor:{sha}"]
      ieng["SA {prefixo}-agente"]
    end
    subgraph AR["no - Artifact Registry"]
      iar["hub, api e agente-professor,<br/>por {sha}"]
    end
    subgraph BAG["no - Cloud Storage -agente-professor"]
      ibag["privado: BNCC, livros e habilidades"]
    end
    subgraph VS["servico gerenciado - Vertex AI Search"]
      ivs["data stores em global:<br/>-habilidades (estruturado) e -acervo (PDFs)"]
    end
    subgraph SM["no - Secret Manager"]
      ism["dez segredos, env do hub e da API"]
    end
  end

  %% caminhos de comunicacao
  LAV ---|"WebRTC (video)"| NAV
  NAV ---|"HTTPS e WSS"| LB
  LB ---|hub| RUNHUB
  LB ---|api| RUNAPI
  LB ---|CDN| GCS
  RUNHUB ---|API_INTERNAL_URL| RUNAPI
  RUNAPI ---|"assets, planos"| GCS
  RUNAPI ---|"publica; push"| PS
  RUNAPI ---|VPC| SQL
  JOB ---|VPC| SQL
  RUNAPI ---|"Gemini e Live"| VX
  RUNAPI ---|"SSE (:streamQuery)"| ENG
  ENG ---|Gemini| VX
  ENG ---|busca| VS
  VS ---|le| BAG
  LAV ---|"HTTPS: /v1/sessions (API -> vendor) e /ajuda/llm (vendor -> LB -> API)"| RUNAPI

  %% o CI implanta cada artefato
  bgcs ==>|baixa| bnav
  GHA ==>|rsync| GCS
  GHA ==>|push| AR
  GHA ==>|"PATCH da imageUri"| aeng

  %% A mesma imagem api:{sha} esta implantada em dois nos, e o bundle esta
  %% armazenado num e executado em outro. O CI ainda executa o job de migracao e
  %% publica o manifesto (PUT /experiencias/:expId), o que o diagrama deixa no
  %% texto. Escala, rede e certificados sao do Terraform, em outro repositorio.
  %% So a API chama o Agent Engine (ADC da SA da api, :query e :streamQuery?alt=sse);
  %% o navegador nunca fala com ele nem com os data stores. O mascote da Ajuda
  %% (ADR-158) e o WSS do navegador a API pelo LB, e da API ao Gemini Live.
  %% Nada do LiveAvatar e implantado pelo CI: o contexto e a configuracao de LLM
  %% do vendor saem de scripts/liveavatar-configurar.mjs (ADR-132). O WSO2 nao
  %% aparece: e o emissor futuro; hoje quem assina os tokens e o hub.
\end{lstlisting}

% ---------------------------------------------------------------------------
\section{Figura~\ref{fig:deploy-contrato} --- Implantação de um contrato}

\begin{lstlisting}[language=mermaid]
C4Deployment
  title Implantacao de um contrato

  Deployment_Node(internet, "Internet", "escola") {
    Container(browser, "Navegador da escola", "", "")
  }

  Deployment_Node(gcp, "Google Cloud - projeto do contrato", "GCP") {
    Container(lb, "Cloud Load Balancing", "", "global, 1 IP, host rules: hub, API, bundles, assets")
    Container(gcs, "Cloud Storage", "", "-exp e -assets atras da CDN; -planos privado")
    Container(hub, "Cloud Run {contrato}-frontend", "hub (Next.js)", "")
    Container(api, "Cloud Run {contrato}-api", "API (NestJS)", "")
    ContainerQueue(ps, "Pub/Sub", "", "telemetria analitica")
    Deployment_Node(vpc, "VPC (so IP privado)", "") {
      ContainerDb(sql, "Cloud SQL {contrato}-pg", "HA, PITR, pooler", "")
    }
    Deployment_Node(ia, "IA gerenciada - Vertex AI", "") {
      Container(agente, "Agent Engine agente-professor", "Python + ADK", "")
      Container(gemini, "Gemini", "", "global; o Live em us-central1")
      Container(busca, "Vertex AI Search", "", "dois data stores em global")
      Container(bag, "Cloud Storage -agente-professor", "", "BNCC e acervo; privado")
    }
  }

  System_Ext(idp, "WSO2 Identity Server", "projeto proprio, varios contratos")
  System_Ext(lav, "LiveAvatar (HeyGen)", "video da Ajuda; fora do Google (ADR-132)")

  Rel(browser, lb, "443")
  Rel(lb, gcs, "bundles, assets (CDN)")
  Rel(lb, hub, "HTTPS")
  Rel(lb, api, "HTTPS")
  Rel(hub, api, "BFF")
  Rel(hub, idp, "OIDC")
  Rel(api, gcs, "assets, planos")
  BiRel(api, ps, "publica / recebe push")
  Rel(api, sql, "5432 - IP PRIVADO")
  Rel(api, gemini, "HTTPS; WSS do Gemini Live (ADR-158)")
  Rel(api, agente, "SSE (ADR-126)")
  Rel(agente, gemini, "modelo")
  Rel(agente, busca, "busca")
  Rel(bag, busca, "le")
  BiRel(api, lav, "sessoes; /ajuda/llm pelo LB")
  Rel(browser, lav, "WebRTC")
\end{lstlisting}

% ---------------------------------------------------------------------------
\section{Figura~\ref{fig:deploy-detalhe} --- Implantação detalhada, em execução}

\begin{lstlisting}[language=mermaid]
flowchart TB
  browser([Navegador]); wso2[[WSO2 IS - projeto proprio]]

  subgraph GCP["Google Cloud - projeto do contrato"]
    subgraph BORDA["Borda global"]
      armor["Cloud Armor (WAF nos backends)"]
      lb["Load Balancer (1 IP, host rules; um certificado por dominio)"]
      cdn["Cloud CDN (bundles e assets)"]
    end
    gcs["Cloud Storage: -exp e -assets publicos atras da CDN;<br/>-planos privado e versionado; -exp-mapas privado;<br/>-agente-professor privado"]
    hub["Cloud Run -frontend (hub; min=1)"]
    api["Cloud Run -api (CPU sempre alocada, min>=1, conc. 80)"]
    job["Cloud Run Job -migrate (prisma migrate deploy)"]
    subgraph VPC["VPC"]
      conn["Serverless VPC Access (conector)"]
      sql[("Cloud SQL -pg (HA, PITR, pooler; sem IP publico)")]
    end
    subgraph IA["IA gerenciada - Vertex AI"]
      eng["Agent Engine agente-professor (us-east1; SA -agente)"]
      gemini["Gemini (global, reserva em us; Live em us-central1)"]
      busca["Vertex AI Search (-habilidades e -acervo; global)"]
    end
    ps["Pub/Sub: topic -telemetria, push sub -telemetria-ingest"]
  end

  browser -->|443| armor --> lb --> cdn --> gcs
  lb --> hub & api
  hub -->|BFF| api
  hub -->|OIDC| wso2
  api -->|-assets, -planos| gcs
  api -->|SSE| eng
  eng --> gemini
  eng -->|busca| busca
  gcs -.->|le -agente-professor| busca
  api -->|privado| conn -->|5432| sql
  job --> conn
  api -->|HTTPS, WSS| gemini
  api <-->|publica; push /telemetry/ingest| ps
\end{lstlisting}

% ---------------------------------------------------------------------------
\section{Figura~\ref{fig:deploy-publicacao} --- O que o CI publica}

\begin{lstlisting}[language=mermaid]
flowchart LR
  gh[[GitHub Actions - deploy.yml; id-token: write]]
  wif["WIF pool github: troca o OIDC pelo token curto da DEPLOY_SA"]
  ar["Artifact Registry {contrato}: imagens da API, do hub e do agente"]
  job["Cloud Run Job -migrate (antes do deploy da API)"]
  run["Cloud Run -api, -frontend (nova revisao + sonda)"]
  eng["Agent Engine (so se apps/agente-professor mudou)"]
  exp["Cloud Storage -exp (bundles), -exp-mapas"]
  api["API: PUT do manifesto (versaoAtiva)"]

  gh -->|OIDC| wif
  wif -->|docker push| ar
  wif -->|jobs execute| job
  wif -->|run deploy| run
  wif -->|PATCH| eng
  wif -->|storage rsync| exp
  gh -->|token RS256 de 5 min| api
  ar -.->|imagem| job & run & eng
\end{lstlisting}

% ---------------------------------------------------------------------------
\section{Figura~\ref{fig:identidade-maquina} --- Identidade de carga de trabalho}

\begin{lstlisting}[language=mermaid]
flowchart LR
  subgraph CI["CI - publicar"]
    gh["GitHub Actions (id-token: write)"] --> tok["token OIDC do GitHub (repo, ref, workflow)"]
    tok --> wif["WIF pool github (valida e troca)"]
    wif --> sadeploy["access token da DEPLOY_SA (curta duracao)"]
  end

  subgraph APP["API - chamar a IA"]
    run["Cloud Run -api (sem env de credencial)"] --> md["metadata server (ADC nivel 3)"]
    md --> saapi["SA {contrato}-api (aiplatform.user; agenteProfessorInvocador)"]
    saapi --> vx["Gemini e Agent Engine - autorizados"]
  end

  subgraph AG["Agente - consultar modelo e acervo"]
    eng["Agent Engine (o agente em execucao)"] --> rt["credencial do runtime (a SA do engine)"]
    rt --> saag["SA {prefixo}-agente (aiplatform.user; discoveryengine.viewer)"]
    saag --> ds["Gemini e os dois data stores - autorizados"]
  end

  subgraph PS["Pub/Sub - entregar a telemetria"]
    sub["push sub -telemetria-ingest"] --> idt["ID token da SA {contrato}-push (oidc_token)"]
    idt --> guard["PubSubPushGuard (emissor Google, aud e email)"]
    guard --> ing["POST /telemetry/ingest - aceito"]
  end

  %% A mesma cadeia de ADC nos tres ambientes:
  %%   1. GOOGLE_APPLICATION_CREDENTIALS
  %%   2. arquivo de `gcloud auth application-default login`
  %%   3. metadata server (Cloud Run)
  %% NUNCA chave de API.
\end{lstlisting}

% ---------------------------------------------------------------------------
\section{Figura~\ref{fig:tenancy-niveis} --- Os três níveis de isolamento}

\begin{lstlisting}[language=mermaid]
flowchart TB
  %% Nivel 1 - fisico: uma stack GCP por contrato
  subgraph A["Contrato A - stack GCP"]
    apiA["API do contrato A"] --> pgA[("PostgreSQL do contrato A")]
    subgraph ESC["Nivel 2 - logico: as escolas dividem a stack"]
      e1["escola sp-001"]; e2["escola sp-002"]; e3["... 100+"]
    end
    tab["Toda tabela de dado de estudante:<br/>escola_id TEXT NOT NULL<br/>ENABLE + FORCE ROW LEVEL SECURITY<br/>POLICY USING (escola_id = current_setting('app.escola_id', true))"]
    e1 -.-> tab; e2 -.-> tab; e3 -.-> tab
    subgraph N3["Nivel 3 - autorizacao: dentro da escola, por papel"]
      p1["estudante: a propria matricula e o progresso"]
      p2["professor: so as turmas em que leciona"]
      p3["gestor_escola: a escola inteira"]
    end
  end
  subgraph B["Contrato B - OUTRA stack"]
    apiB["API do contrato B"] --> pgB[("PostgreSQL do contrato B")]
  end
\end{lstlisting}

% ---------------------------------------------------------------------------
\section{Figura~\ref{fig:tenancy-fluxo} --- Propagação da escola}

\begin{lstlisting}[language=mermaid]
flowchart LR
  req["Requisicao HTTP (Authorization: Bearer)"] --> guard["OidcGuard (JWKS; injeta req.auth)"]
  guard --> roles["RolesGuard (le req.auth.perfil)"]
  roles --> int["TenancyInterceptor (definirTenancyNoContexto)"]
  int --> ctrl["controller"] --> svc["service: comTenancy(fn)"]
  svc --> tx["Transacao: set_config + query"]
  svc -->|sem escola no contexto| fail["LANCA"]

  %% Guards rodam ANTES de interceptors, entao req.auth ja existe.
  %% Rota sem OidcGuard nao abre contexto — e comTenancy sem escola falha alto.
  %% A RLS e rede de seguranca, nao a unica tranca: os services continuam
  %% filtrando escola_id explicitamente em toda query.
\end{lstlisting}

% ---------------------------------------------------------------------------
\section{Figura~\ref{fig:er-plataforma} --- Modelo de dados da plataforma}

\begin{lstlisting}[language=mermaid]
erDiagram
  escolas ||--o{ turmas : tem
  escolas ||--o{ users  : "escola_id; regional_id do gestor regional"
  turmas  ||--o{ users  : "turma_id (cache da principal)"
  escolas ||--o{ agente_professor_conversas : tem
  users   ||--o{ agente_professor_conversas : "dono da sessao"
  experiencias ||..o{ progresso : "exp_id (sem FK)"
  telemetria_eventos ||..o{ particoes : "ATTACH por mes"

  escolas {
    text id PK "MESMO valor do claim escola_id e do set_config"
    text nome
    text codigo_inep
    text tipo "rede|regional|escola - so a escola e tenant"
    text pai_id FK "escolas: rede > regional > escola"
    timestamptz created_at
  }
  turmas {
    uuid id PK
    text escola_id FK
    uuid periodo_letivo_id FK
    text nome "UNIQUE (escola_id, periodo_letivo_id, nome)"
    text serie "1EF..9EF | 1EM..3EM"
  }
  users {
    uuid id PK
    text role "estudante|professor|gestor_escola|gestor_regional|gestor_contrato|admin"
    text escola_id FK "NULL so p/ gestor_regional, gestor_contrato, admin (CHECK)"
    text regional_id FK "escolas, p/ gestor_regional"
    uuid turma_id FK "CACHE da matricula principal"
    text username "UNIQUE (escola_id, username)"
    text display_name
    text password_hash
    text cpf
  }
  experiencias {
    text exp_id PK
    text nome
    text versao_ativa "rollback = UPDATE deste campo"
    text produto
    smallint unidade
    text_array series "VAZIO = TODAS as series"
    text_array capabilities
    jsonb viewport "presente => casca de aparelho"
    text publico "estudante|professor"
    text_array permissoes
  }
  progresso {
    text escola_id PK
    text estudante_id PK
    text exp_id PK
    jsonb dados
    timestamptz atualizado_em
  }
  ia_uso_diario {
    text escola_id PK
    text usuario_id PK
    text dia PK "YYYY-MM-DD"
    int tokens_usados
  }
  agente_professor_conversas {
    uuid id PK
    text escola_id FK
    uuid usuario_id FK
    text sessao_id "UNIQUE: uma sessao tem um dono so"
    text titulo "CACHE do titulo_conversa do agente"
    timestamptz criada_em
    timestamptz atualizada_em
  }
  telemetria_eventos {
    uuid event_id PK "PK COMPOSTA: o Postgres exige a chave de particao"
    timestamptz ts PK
    text escola_id
    text estudante_id
    text exp_id
    text exp_version
    uuid sessao_id
    text verbo
    text objeto
    jsonb resultado
    jsonb contexto
  }
  particoes {
    table telemetria_eventos_YYYY_MM "uma por mes, fora do Prisma"
    table telemetria_eventos_default "rede de seguranca"
  }
\end{lstlisting}

% ---------------------------------------------------------------------------
\section{Figura~\ref{fig:er-enturmamento} --- Enturmamento e sala de aula}

\begin{lstlisting}[language=mermaid]
erDiagram
  periodos_letivos ||--o{ turmas     : "a turma vive dentro de um ano letivo"
  turmas           ||--o{ matriculas : enturma
  users            ||--o{ matriculas : "pessoa x turma, datada"
  turmas           ||--o{ liberacoes : "a missao DENTRO da sala"
  users            ||--o{ liberacoes : "quem liberou"

  periodos_letivos {
    uuid id PK
    text tipo "ano_letivo|semestre|bimestre|trimestre"
    text titulo "UNIQUE (tipo, titulo)"
    smallint ano
    date inicio
    date fim
    uuid pai_id FK "ano letivo -> bimestre; NULL = raiz"
    text fonte_id
  }
  turmas {
    uuid id PK
    text escola_id FK
    uuid periodo_letivo_id FK "NOT NULL: turma sem periodo e users.turma_id de novo"
    text nome "UNIQUE (escola_id, periodo_letivo_id, nome)"
    text serie "1EF..9EF | 1EM..3EM"
    text tipo "regular|componente (classType do OneRoster)"
    text turno
    text codigo
    text fonte_id
  }
  matriculas {
    uuid id PK "RLS por escola_id"
    text escola_id FK
    uuid turma_id FK
    uuid usuario_id FK
    text papel "estudante|professor — gestor NAO entra"
    boolean principal "alimenta o cache users.turma_id"
    date inicio
    date fim "sair e fim, nunca DELETE"
    text fonte_id
  }
  liberacoes {
    uuid id PK "RLS por escola_id"
    text escola_id FK
    uuid turma_id FK
    text exp_id "SEM FK: o catalogo experiencias e escrito pelo CI"
    text estado "liberada|encerrada — nao existe bloqueada"
    uuid liberada_por FK
    timestamptz liberada_em
    timestamptz encerrada_em "CHECK: casa com estado"
  }
  users {
    uuid id PK
    text matricula_externa "o RA da rede; indice unico parcial por escola"
    uuid turma_id FK "CACHE da matricula principal — a verdade e matriculas"
    text fonte_id
    text email
    boolean ativo
    boolean senha_provisoria
  }

  %% A ausencia de linha em liberacoes E o bloqueio: guardar bloqueio como linha
  %% obrigaria a criar uma por turma x experiencia so para o mundo nascer fechado.
  %% liberacoes e PRE-CONDICAO do relogio da exp02em1b e dos gates da exp06ef2b, que
  %% dizem "a aula comecou" e continuam morando na tabela DA experiencia.
\end{lstlisting}

% ---------------------------------------------------------------------------
\section{Figuras~\ref{fig:er-exp08ef1ba} e~\ref{fig:er-exp08ef1bb} --- Modelo da exp08ef1b}

\begin{lstlisting}[language=mermaid]
erDiagram
  exp08ef1b_activities ||--o{ exp08ef1b_authors : tem
  exp08ef1b_authors    ||--o{ exp08ef1b_posts   : escreve
  exp08ef1b_posts      ||--o{ exp08ef1b_post_comments : tem
  exp08ef1b_activities ||--o{ exp08ef1b_generation_jobs : enfileira
  exp08ef1b_activities ||--o{ exp08ef1b_runs : hospeda
  exp08ef1b_runs ||--o{ exp08ef1b_events : registra
  exp08ef1b_runs ||--o{ exp08ef1b_student_posts : publica
  exp08ef1b_runs ||--|| exp08ef1b_conclusoes : conclui
  exp08ef1b_runs ||--|| exp08ef1b_reflexoes : reflete
  exp08ef1b_runs ||--|| exp08ef1b_dossiers : gera
  exp08ef1b_runs ||--o{ exp08ef1b_comentarios_da_sessao : ancora
  exp08ef1b_activities ||--o{ exp08ef1b_app_70_perfis_importados : "largada"
  exp08ef1b_activities ||--o{ exp08ef1b_posts : contem
  exp08ef1b_authors    |o--o{ exp08ef1b_post_comments : "author_id; NULL = handle avulso"
  exp08ef1b_posts      |o--o{ exp08ef1b_events : "post_id"
  exp08ef1b_posts      |o--o{ exp08ef1b_comentarios_da_sessao : "post_id (XOR)"
  exp08ef1b_student_posts |o--o{ exp08ef1b_comentarios_da_sessao : "student_post_id (XOR)"
  exp08ef1b_runs ||..|| exp08ef1b_app_70_perfis_importados : "mesmo estudante e atividade (sem FK)"

  exp08ef1b_activities {
    uuid id PK
    text escola_id FK
    uuid professor_id FK
    text pergunta_chave
    text pergunta_origem "ia|professor - auditoria de autoria"
    text etapa "relogio COLETIVO da turma"
    boolean ia_ao_vivo "desligavel sem deploy"
    jsonb antagonismo
    jsonb engine_config
    text status
  }
  exp08ef1b_posts {
    uuid id PK
    text lado "VARIAVEL INDEPENDENTE - sanitizada antes de sair"
    text perfil "idem"
    smallint intensidade "idem"
    text image_status "pendente|ok|fallback"
    boolean aprovado
  }
  exp08ef1b_runs {
    uuid id PK
    text lado_inicial
    text lado_origem "perfil|sorteio|redistribuido - auditoria"
    jsonb engine_state "estado do motor; unidade de LOCK"
    text nome_no_app "a conta e do APARELHO, nao da crianca"
  }
  exp08ef1b_events {
    bigint id PK
    text tipo "canal OPERACIONAL: sincrono e transacional"
    jsonb payload "{dwellMs, dwellNorm}"
  }
  exp08ef1b_comentarios_da_sessao {
    uuid post_id FK "XOR com student_post_id (CHECK na migration)"
    uuid student_post_id FK
    boolean de_ia "para o professor distinguir na auditoria"
  }
  exp08ef1b_app_70_perfis_importados {
    jsonb escores "4 quadrantes, 0-1"
    text dominante
    text origem "jogo|professor"
    text lado_forcado "PRESERVADO no upsert"
  }
\end{lstlisting}

% ---------------------------------------------------------------------------
\section{Figuras~\ref{fig:er-exp04ef1b-app70} e~\ref{fig:er-app90-91} --- Demais experiências}

\begin{lstlisting}[language=mermaid]
erDiagram
  app_70_perfis  ||..|| app_70_bolhas : "mesmo estudante"
  app_70_sessoes ||--o{ app_70_perfis : "sessao_id"
  app_90_planos  ||..o{ app_90_referencias : "monta o prompt"
  app_91_ciclos  ||--o{ app_91_checkins : "7 check-ins"

  backup_exp04ef1b_sessoes {
    text escola_id PK
    uuid estudante_id PK
    text_array salvos "CHECK cardinality <= 3"
    boolean ranking_desbloqueado "LATCH: nao volta"
    boolean credibilidade_adicionada "congela os salvos"
  }
  app_70_sessoes {
    uuid id PK "EMITIDO pela API - o cliente colidia"
    uuid turma_id "SEM FK, desnormalizado: turma NO MOMENTO da partida"
    text bolha "A|B|C - do SERVIDOR"
    smallint cartoes_servidos "!= apareceram"
    timestamptz aberta_em "relogio do SERVIDOR"
    jsonb payload "bruto - unica fonte do reprocessamento"
    jsonb leitura_mensagens
  }
  app_70_perfis {
    smallint versao_mapeamento "sobe => nasce linha nova"
    jsonb valores "null NAO e zero"
    jsonb graus "scored|indicative|zero|none"
    float abertura_vs_conservacao "NULL se polo nao observado"
    boolean amostra_suficiente "falso => o painel avisa que NAO SABE"
  }
  app_70_bolhas {
    text escola_id PK
    uuid estudante_id PK "por ESTUDANTE, nao por partida"
    text bolha
  }
  app_90_planos {
    text status "pendente|gerando|pronto|erro"
    text caminho_gerado_md "IMUTAVEL, bucket privado"
    text caminho_html "o que o editor salvou, sanitizado"
  }
  app_90_referencias {
    text disciplina PK
    text serie PK "NULL = todas"
    text nome PK
    text texto "catalogo do CONTRATO - sem escola_id, FORA da RLS"
  }
  app_91_ciclos {
    uuid id PK
    int_array ordem_itens "permutacao de 1..21"
    int depressao "somas x2, ao concluir - DADO SENSIVEL"
  }
  app_91_checkins {
    smallint numero "1..7"
    text dia "UNIQUE (escola_id, professor_id, dia)"
    text texto_livre "NUNCA em log nem telemetria"
    text origem_texto "ia|fallback|seguranca"
    boolean sinal_de_apoio
  }
\end{lstlisting}

% ---------------------------------------------------------------------------
\section{Figura~\ref{fig:seq-handshake} --- \emph{Handshake} hub-experiência}

\begin{lstlisting}[language=mermaid]
sequenceDiagram
  autonumber
  participant E as Experiencia (@edu/sdk no iframe)
  participant F as ExpFrame (React, no hub)
  participant R as Rota do hub (servidor)
  participant A as API (NestJS)

  R->>A: GET /experiencias/:expId
  A-->>R: manifesto { versaoAtiva, viewport, permissoes }
  R->>R: assina o token (RS256, 15 min)
  R->>F: props: src, token, apiUrl
  F->>E: renderiza o iframe com a URL versionada
  E->>E: connectToHub({ hubOrigin })
  loop a cada 400ms, ate o hub responder
    E->>F: exp:ready { protocolVersion: 2 }
  end
  F->>F: valida ev.origin e a versao MAIOR
  F-->>E: hub:init { token, estudante, apiUrl, config }
  E->>E: valida com Zod; resolve a Promise
  E->>A: chamadas de dominio, /telemetry/batch, /ai/gerar
  E->>F: exp:token-expirando
  F->>R: GET /api/exp-token
  R-->>F: { token }
  F-->>E: hub:token { token }

  Note over E,F: Com assets em cache (F5), a experiencia fica pronta ANTES<br/>de o ExpFrame hidratar. Um disparo unico se perderia.<br/>hub:init repetido e evento normal, nao ruido.
\end{lstlisting}

% ---------------------------------------------------------------------------
\section{Figura~\ref{fig:fluxo-embed} --- Resolução da URL do \emph{embed}}

\begin{lstlisting}[language=mermaid]
flowchart TB
  ini([inicio]) --> le["le em paralelo na api, com o token recem-assinado:<br/>GET /experiencias/:expId (manifesto) e<br/>GET /recursos/efetivos (versao fixada na escola)"]
  le --> q1{tem manifesto?}
  q1 -->|sim| vp1["viewport e permissoes vem do MANIFESTO"]
  q1 -->|nao| vp2["le {baseInterna}/experiencia.json<br/>(EXP_INTERNAL_HOST aplicado a base)"]
  vp1 --> q2{producao e com manifesto?}
  vp2 --> q2
  q2 -->|sim| com["{cdn}/{expId}/{versao}/index.html?hubOrigin=...<br/>versao: a fixada no escopo, senao a versaoAtiva"]
  q2 -->|nao| sem["{base}/?hubOrigin=...<br/>base: em dev, o src de /dev-embed ou a porta de<br/>lib/experiencias-dev.ts; em producao, o CDN"]
  com --> frame["ExpFrame: sandbox + allow,<br/>dentro de DispositivoFrame se houver viewport"]
  sem --> frame
  frame --> fim([fim])

  %% A leitura do experiencia.json NAO e so de dev: stack com a tabela vazia
  %% servia a exp08ef1b esticada na janela em vez de dentro do aparelho.
\end{lstlisting}

% ---------------------------------------------------------------------------
\section{Figura~\ref{fig:contrato-exps} --- Contrato entre duas experiências}

\begin{lstlisting}[language=mermaid]
flowchart LR
  prof(["Professor"]) -->|"POST /atividades/:id/perfis/da-vila"| mv
  v["app-70 (Vila do Lume)"] -->|payload bruto| mv["vila/quadrantes.ts (modulo vila da api): colapso 10 -> 4"]
  mv -->|"grava (upsert)"| t["exp08ef1b_app_70_perfis_importados (a tabela de contrato)"]
  t -->|le| mi["perfis/largada.ts (modulo perfis da api): decide o lado do feed"]
  mi -->|lado inicial| i["exp08ef1b (Instanooria)"]

  %% Duas travas que CUSTAM estudantes, de proposito:
  %%  1. valor nao observado (indicative|none) fica FORA da media — nunca zero
  %%  2. quadrante sem nenhum membro observado INVALIDA o perfil inteiro
  %% Sem perfil, o Instanooria cai no sorteio balanceado, que e honesto sobre
  %% nao saber. E `lado_forcado` e PRESERVADO no upsert.
\end{lstlisting}

% ---------------------------------------------------------------------------
\section{Figura~\ref{fig:telemetria-analitica} --- Canal analítico}

\begin{lstlisting}[language=mermaid]
flowchart LR
  exp["Experiencia: track(verbo, objeto, dados)"] --> sdk["TelemetryClient<br/>buffer em memoria; lote de 50 ou 10s"]
  sdk -->|"<= 50 eventos"| batch["POST /telemetry/batch<br/>OidcGuard; identidade vem do TOKEN"]
  batch --> topic[["Pub/Sub topic -telemetria"]]
  sdk -->|sendBeacon| beacon["POST /telemetry/beacon<br/>pagehide; token NO CORPO; responde 204 SEMPRE"]
  beacon -->|publica no mesmo topic| topic
  topic -->|push autenticado| ingest["POST /telemetry/ingest<br/>PubSubPushGuard: ID token do Google, aud, email"]
  ingest --> pg[("telemetria_eventos<br/>INSERT multi-linha<br/>ON CONFLICT (event_id, ts) DO NOTHING<br/>particionada por mes")]

  %% Onde ha perda possivel, e SO ali: entre o track() e o flush.
  %% Do Pub/Sub em diante a entrega e DURAVEL — e por isso um buffer em Redis
  %% nao o substituiria: reintroduziria a perda que ele eliminou.
\end{lstlisting}

% ---------------------------------------------------------------------------
\section{Figura~\ref{fig:camada-ia} --- Camada de IA}

\begin{lstlisting}[language=mermaid]
C4Component
  title Camada de IA - src/ai e a UNICA porta

  Container_Boundary(api, "API - apps/api/src") {
    Component(ger, "geracao/, motor/", "", "exp08ef1b")
    Component(dom, "modulos de dominio", "", "planos/, bemestar/ e os demais")
    Component(sdk, "POST /ai/gerar", "", "a rota da SDK")
    Component(age, "agente-professor/", "", "o chat do professor")
    Component(aju, "ajuda/", "", "mascote e avatar")

    Container_Boundary(ai, "src/ai/ - a UNICA porta para a IA") {
      Component(q,   "QuotaService", "", "temSaldo ANTES, consumir DEPOIS; teto em ia_quotas, soma em ia_uso_diario")
      Component(pr,  "AiProvider", "", "generate, generateJson, conversar, generateImage, generateSpeech; embrulhado pelo MedidorDeProvider")
      Component(pag, "AgenteProfessorPort", "", "")
      Component(pao, "VozAoVivoPort / AvatarAoVivoPort", "", "")
      Component(vx,  "vertex.provider.ts", "", "default; ADC + REST")
      Component(mk,  "mock.provider.ts", "", "so testes")
      Component(pp,  "prompts.ts", "", "prompt de sistema; nao sai daqui")
      Component(aec, "agent-engine.client.ts", "", "SSE, ADC")
      Component(liv, "gemini-live.client.ts / liveavatar.client.ts", "", "")
    }
  }
  System_Ext(vertex, "Vertex AI - Gemini", "global, reserva em us")
  System_Ext(eng, "Agent Engine", "o agente do professor")
  System_Ext(glive, "Gemini Live", "us-central1")
  System_Ext(lav, "LiveAvatar", "HeyGen")

  Rel(ger, pr, "sem quota (ADR-74)")
  Rel(dom, q, "")
  Rel(sdk, q, "")
  Rel(age, q, "")
  Rel(q, pr, "so com saldo")
  Rel(q, pag, "")
  Rel(aju, pao, "minutos por dia (ADR-132)")
  Rel(pr, vx, "AI_PROVIDER=vertex")
  Rel(pr, mk, "so testes")
  Rel(pr, pp, "usa")
  Rel(pag, aec, "")
  Rel(pao, liv, "")
  Rel(vx, vertex, "HTTPS")
  Rel(aec, eng, ":streamQuery")
  Rel(liv, glive, "WSS")
  BiRel(liv, lav, "")
  %% Todas as portas registram o custo em ia_consumo_diario (ADR-161).
\end{lstlisting}

% ---------------------------------------------------------------------------
\section{Figura~\ref{fig:guardrails-bemestar} --- \emph{Guardrails} do bem-estar}

\begin{lstlisting}[language=mermaid]
flowchart TB
  ini([inicio]) --> abre["GET /bemestar/hoje - 3 itens da permutacao do ciclo"]
  abre --> sinal{"sinal de apoio?<br/>item 17 ou 21 >= 2, ou termo de risco"}
  sinal -->|sim| fixo["texto FIXO de apoio (CVV 188, SAMU 192, CAPS/UBS)<br/>a IA NAO e chamada<br/>origem_texto = 'seguranca'"]
  sinal -->|nao| saldo{tem saldo e provider?}
  saldo -->|nao| fall["texto FIXO do catalogo<br/>origem_texto = 'fallback'"]
  saldo -->|sim| ia["IA gera abertura, pergunta aberta e acolhimento<br/>recebe SO: dia da semana, mes, no do check-in,<br/>texto livre delimitado"]
  ia -->|falha| fall
  ia --> val{validarTextoDaIa?}
  val -->|barrado| fall
  val -->|passou| ok["origem_texto = 'ia'"]
  fixo --> fim([fim]); fall --> fim; ok --> fim

  %% Os 21 itens do DASS-21 sao FIXOS e VERBATIM. A IA nunca reescreve um item:
  %% reformular quebra a validacao psicometrica e invalida as faixas de corte.
  %% O check-in NUNCA deixa de abrir.
\end{lstlisting}

% ---------------------------------------------------------------------------
\section{Figura~\ref{fig:graus} --- Os quatro graus de medida}

\begin{lstlisting}[language=mermaid]
classDiagram
  class scored {
    >= 2 sinais independentes
    numero: number
    entra na media
    E PUBLICADO com numero
  }
  class indicative {
    exatamente 1 sinal
    numero: null
    descrito em TEXTO
    fora dos eixos e das medias
  }
  class zero {
    teve a chance e NAO foi
    numero: null
    ENTRA na media, como zero de verdade
  }
  class none {
    a oportunidade NAO chegou
    numero: null
    SAI da conta - nao e zero
  }
  note for indicative "Por que o tipo, e nao a tela: no app de origem a decisao vivia\nna interface, e uma tela nova podia imprimir 'Poder 100' de um sinal so."
  note for none "A distincao mais importante do modelo e entre zero e none.\nTratar 'nao teve chance' como 'escolheu nao' faz uma crianca\nparecer ter recusado algo que nunca lhe foi oferecido."
\end{lstlisting}

% ---------------------------------------------------------------------------
\section{Figura~\ref{fig:app92-caminho} --- app-92: o caminho do planejamento}

\begin{lstlisting}[language=mermaid]
flowchart TB
  subgraph fora["Fora do produto, uma vez"]
    pdf[PDFs homologados do MEC] --> ext[scripts/extrair-bncc.mjs]
    ext --> dados[("api/src/curriculo/dados/*.ts<br/>1.304 EF - 183 EM - 129 Computacao")]
  end

  ini((inicio))
  ini --> reg["1. Regencia<br/>componente, serie, ano letivo, grade de horarios,<br/>minutos por aula, divisao do ano<br/>POST /planejamento/configuracoes"]
  reg --> dist["2. Distribuicao<br/>o professor poe cada habilidade num periodo - DECISAO HUMANA<br/>PUT .../planos-anuais/:id/distribuicao"]
  dist --> anual["3. Plano anual (IA)<br/>a prosa pedagogica do ano, gerada e depois reescrita pelo professor<br/>POST .../planos-anuais/:id/gerar"]
  anual --> aulas["4. Aulas, semana a semana (IA)<br/>40 semanas, uma por vez. O esqueleto e do calendario<br/>POST .../periodos/:ordem/gerar"]
  aulas --> exp["Exportacao<br/>montada no NAVEGADOR, sem rota e sem modelo:<br/>.doc (HTML com os namespaces do Office) e impressao"]
  exp --> fim(((fim)))

  dados --> dist

  reg --> cfg[("app_92_configuracoes<br/>e, se a rede ainda nao tinha, periodos_letivos - catalogo do CONTRATO")]
  dist --> jd[("app_92_planos_anuais.distribuicao (JSONB)<br/>o insumo do gerador, e o que sobrevive a reescrever o documento")]
  anual --> doc[("bucket privado: gerado.md (imutavel) + atual.html (sanitizado)")]
  aulas --> ja[("app_92_planos_de_periodo.aulas (JSONB)<br/>dado estruturado, nao prosa")]
  exp --> rec["tres recortes, so das semanas pronto:<br/>esta semana - esta e a proxima -<br/>o periodo de avaliacao inteiro (bimestre, trimestre ou semestre)"]

  %% Onde a IA NAO entra: a extracao da Base; a consulta de habilidades
  %% (exata por componente e ano, sem busca vetorial); a distribuicao; o
  %% esqueleto das aulas; a costura; cumprida e observacaoDoProfessor; e a
  %% exportacao.
  %% Os dois pontos de IA passam por temSaldo ANTES e consumir depois - e o
  %% consumo conta mesmo quando a geracao e rejeitada: o token foi gasto.
\end{lstlisting}

% ---------------------------------------------------------------------------
\section{Figura~\ref{fig:app92-semana} --- app-92: como uma semana nasce}

\begin{lstlisting}[language=mermaid]
flowchart TB
  cal["Calendario (calendario.ts)<br/>numero, data, dia da semana e horario das aulas DESTA semana,<br/>da grade real do professor. Feriado REMOVE a aula, nao a desloca"]
  rec["O recorte da distribuicao<br/>so os codigos do periodo de avaliacao em que esta semana cai"]
  reg["A regencia<br/>componente, serie e minutos por aula"]

  cal --> brief
  rec --> brief
  reg --> brief

  brief["montarBriefingDePeriodo<br/>um briefing com as aulas que EXISTEM e as habilidades que PODEM entrar"]
  brief --> gen["provider.generateJson (IA)<br/>schema de resposta, modelo longo, temperatura 0.5<br/>devolve SO o conteudo pedagogico das aulas"]
  gen --> cost["costurarAulas - a conciliacao entre o que o modelo devolveu<br/>e as aulas que existem de verdade"]
  cost --> out[("app_92_planos_de_periodo.aulas<br/>JSONB estruturado, status = pronto")]
  out --> edita["O professor edita - PUT /planejamento/periodos/:id/aulas<br/>cumprida e a observacao da aula sao registro DELE, nunca do modelo"]

  %% As tres regras da costura:
  %% 1. O esqueleto e do CALENDARIO. Oito aulas para uma semana de tres: a
  %%    sobra e descartada. Uma para tres: as outras duas nascem em branco.
  %% 2. Codigo fora do recorte e REMOVIDO, nao corrigido.
  %% 3. Campo ausente vira string VAZIA, nunca undefined.
  %% As tres barreiras contra um codigo inventado: o prompt (instrucao, nao
  %% garantia), o PUT .../distribuicao (valida o que o PROFESSOR escolheu) e a
  %% costura (remove o que o MODELO escreveu depois). Nenhuma sozinha basta.
\end{lstlisting}

% ---------------------------------------------------------------------------
\section{Figuras~\ref{fig:mapa-rotas} e~\ref{fig:mapa-rotas-exps} --- Mapa de rotas da API}

\begin{lstlisting}[language=mermaid]
%% Figura 13.1 - plataforma e exp08ef1b
flowchart TB
  subgraph PLAT["Plataforma - serve a todos os papeis"]
    au["/auth/login - sem papel"]; he["/healthz, /readyz - PUBLICO"]
    ex["/experiencias - GET qualquer; PUT gestor de contrato, admin"]
    pg["/progresso/:expId - estudante"]; te["/telemetry/* - ver capitulo 10"]
    us["/usuarios/* - gestor, admin"]; ai["/ai/gerar - autenticado + quota"]
    an["/analise/* - GESTOR DE CONTRATO, admin"]
    ro["/rostering/* - gestor, admin"]
    tu["/turmas/* - prof. (as dele, por MATRICULA), gestor"]
    mt["/matriculas/:id - gestor, admin"]
    pll["/periodos-letivos - GET prof.; POST gestor"]
    pp["/perfis-recebidos - gestor, admin"]
    rec["/recursos/* - o interruptor da REDE (ADR-87)"]
    sal["/salas/* - prof. libera, gestor le; /minhas: estudante (ADR-69)"]
    eq["/equipes/* - prof. monta, gestor le; /minha: estudante (ADR-111)"]
    ev["/evidencias/* - prof. confirma e ajusta, gestor le; estudante: so a pergunta e a resposta (ADR-112)"]
    agp["/agente-professor/* - SO PROFESSOR, posse; resposta em SSE (ADR-126)"]
    pod["/podcast/:expId/* - prof., gestor; nenhuma de estudante (ADR-150)"]
    aj["/ajuda/* - prof. e gestao; teto em minutos por dia (ADR-132)"]
    obs["/observabilidade/* - autenticado; /publico: sem sessao"]
    prs["/prospeccao/* - /publico: visitante; cadastros: ADMIN (ADR-107)"]
    esc["/escolas/* e /lotacoes - gestao; criar escola: gestor de contrato"]
    est["/estudantes/:id/* - gestao: transferencia e portabilidade"]
  end
  subgraph E01["exp08ef1b - Instanooria"]
    atv["/atividades/* - professor"]; rn["/runs/* - estudante (/eventos: professor tb.)"]
    po["/posts/:id/comentarios - estudante"]; pf["/perfis/importar - chave de servico"]
  end
\end{lstlisting}

\begin{lstlisting}[language=mermaid]
%% Figura 13.2 - demais experiencias: uma linha por experiencia,
%% colunas = estudante | professor e gestao | guia, agregado e servico
flowchart TB
  subgraph UT["backup-exp04ef1b - UnooTube (arquivada, ADR-116)"]
    ut["sem rotas: fora do app.module"]; utt["era /unootube/turmas/*"]
    utp["era /unootube/professor/*"]
  end
  subgraph VI["app-70 - Vila do Lume"]
    vi["/vila/sessoes - estudante"]; vp["/vila/professor/* - PROFESSOR"]
    va["/vila/agregado/* - prof, gestor"]
  end
  subgraph PL["app-90 - Plano de aula"]
    pl["/planos/* - professor"]; pref["/planos/referencias - chave de servico"]
  end
  subgraph BE["app-91 - Bem-estar"]
    be["/bemestar/* - SO PROFESSOR, com posse"]
  end
  subgraph PE["app-92 - Plano de ensino"]
    pe["/planejamento/* - PROFESSOR"]
    pec["/planejamento/curriculo/* - prof, gestor - a Base"]
  end
  subgraph UG["exp02em1b - unooGPT"]
    ug["/unoogpt/* - estudante"]
    ugp["/unoogpt/professor/* - SO PROFESSOR, transcricao"]
    uga["/unoogpt/agregado/* - gestor, admin - sem estudanteId"]
  end
  subgraph RB["exp06ef2b - O Rastro do Boato"]
    rb["/rastro-boato/* - estudante"]
    rbt["/rastro-boato/turmas/* - prof, gestor, admin - o gate da aula"]
    rbp["/rastro-boato/professor/* - prof, gestor - material do guia"]
  end
  subgraph MIA["exp08ef4b - O Manifesto da IA"]
    mia["/manifesto-ia/* - estudante"]
    miat["/manifesto-ia/turmas/* - prof, gestor - o debate E o conteudo da turma"]
  end
  subgraph FC["exp01ef1b - O Que a Foto Conta"]
    fc["/foto-conta/* - estudante"]
    fct["/foto-conta/turmas/* - prof, gestor - o latch da Roda"]
  end
  subgraph JE["exp04ef1b - Jornal da Escola (ADR-117)"]
    je["/jornal-escola/* - estudante - toda a evidência por ação"]
    jet["/jornal-escola/turmas/* - prof, gestor - o pacote, os dois portoes e o painel"]
    jep["/jornal-escola/professor/* - prof, gestor - material do guia"]
  end
  subgraph QM["exp05ef1b - Quem mandou isso?"]
    qm["/quem-mandou-isso/* - estudante - e as perguntas em /evidencias"]
    qmt["/quem-mandou-isso/turmas/* - prof, gestor - o pacote E o relogio da turma"]
    qmp["/quem-mandou-isso/professor/* - prof, gestor - material do guia"]
  end
  subgraph JU["exp06ef1b - O Juri Estadual de Projetos (ADR-137)"]
    ju["/juri-estadual/* - estudante - e as enquetes em /evidencias"]
    jut["/juri-estadual/turmas/* - prof, gestor - a mostra, a Galeria Oculta e a lousa"]
    jup["/juri-estadual/professor/* - prof, gestor - material do guia (so o esqueleto)"]
  end
  subgraph SE["exp07ef1b - A Selecao"]
    sel["/a-selecao/* - estudante"]
    selt["/a-selecao/turmas/* - prof, gestor - as coberturas, as imagens e o portao"]
    selp["/a-selecao/professor/* - prof, gestor - material do guia"]
  end

  %% Duas areas por AUSENCIA deliberada: bem-estar nao tem rota de gestor nem
  %% de agregado; e no unoogpt o controller do gestor nao tem NENHUMA rota que
  %% aceite estudanteId, o que torna estrutural (nao disciplinar) o limite do que
  %% ele alcanca.
\end{lstlisting}

% ---------------------------------------------------------------------------
\section{Figura~\ref{fig:linha-revelacao} --- A linha da revelação}

\begin{lstlisting}[language=mermaid]
flowchart LR
  subgraph CLONE["Dentro do iframe, ANTES da revelacao - @edu/ui/ds-unootube"]
    feed["feed de videos"]; cromo["cromo, busca, menu"]
    play["playlists"]; canal["canal do estudante"]; toast["toast: 'Video salvo!'"]
  end
  subgraph PLATF["DEPOIS da revelacao, no MESMO iframe - @edu/ui/ds"]
    rank["ranking"]; d1["dialogo 1 da missao"]; d2["dialogo 2 da missao"]; fim["fechamento"]
  end
  CLONE ==>|REVELACAO| PLATF

  %% A linha e a REVELACAO, nao o arquivo. O toast NAO cruza: ele fala de
  %% dentro da ficcao. O criterio nao e "e depois no tempo?", e "quem fala?".
  %%
  %% CONTRAEXEMPLO, exp02em1b (UnooriaAI): nao ha revelacao no software. O vies de
  %% cada fonte e o material do debate do professor, e nenhuma tela o entrega.
  %% Logo a exp02em1b e 100% da secao do clone (@edu/ui/ds-unooriaai) e nunca troca
  %% para a da plataforma. A marca institucional entra so como chrome do HUB,
  %% desenhada ACIMA do iframe.
\end{lstlisting}

% ---------------------------------------------------------------------------
\section{Figura~\ref{fig:pipeline-deploy} --- \emph{Pipeline} de entrega}

\begin{lstlisting}[language=mermaid]
flowchart TB
  push["push para dev ou hml (merge de um PR)"]
  push --> ver["PORTAO verificar: ci.yml (typecheck, lint, test, build)"]
  push --> diff["job mudancas: git diff HEAD^ HEAD"]
  join{{"needs: [verificar, mudancas]"}}
  ver --> join
  diff --> join

  subgraph api["job deploy-api - a ORDEM importa"]
    img["imagem da api (camada de deps com cache)"] --> mig["Cloud Run Job -migrate (na VPC, --wait)"]
    mig --> dep["gcloud run deploy (revisao nova da api)"]
    dep --> sapi["sonda pos-deploy: /readyz, ate 2 min"]
  end

  subgraph hub["job deploy-hub - em PARALELO"]
    imgh["imagem do hub (camada de deps com cache)"] --> deph["gcloud run deploy (revisao nova do hub)"]
    deph --> shub["sonda pos-deploy: /api/healthz e JWKS"]
  end

  subgraph agente["job deploy-agente-professor - em PARALELO"]
    imga["imagem do agente (uv sync --locked)"] --> pat["PATCH do Agent Engine (so a imageUri; espera a operacao)"]
  end

  subgraph exps["job publicar-exps (sobe, inerte) -> job ativar-exps"]
    build["turbo build (filtrado, ou tudo se forcar)"]
    build --> rsync["gcloud storage rsync (caminho com a versao, sem .map)"]
    rsync --> cc["Cache-Control: max-age=60 so no index.html"]
    cc --> man["ativar-exps: PUT do manifesto (versaoAtiva)"]
  end

  join --> img
  join --> imgh
  join --> imga
  join --> build
  sapi --> man
  shub --> man
  man --> smoke["job smoke: e2e publico/01 contra o ambiente"]

  %% ORDEM 1: migracao ANTES do deploy - o schema novo tem de estar no banco
  %%          antes de a revisao nova receber trafego. --wait faz falhar.
  %% ORDEM 2: SUBIR o bundle nao espera nada alem do portao (o caminho carrega
  %%          a versao, o upload e inerte); ATIVAR espera api e hub, porque o
  %%          PUT do manifesto e validado contra o JWKS que o HUB serve. Deploys
  %%          PULADOS nao bloqueiam (!cancelled() + result != 'failure').
  %% AGENTE: sem engine ainda (vars.AGENTE_PROFESSOR_ENGINE vazia), o job so
  %%          publica a imagem e avisa; o engine nasce no apply 2 do Terraform.
  %%          O smoke nao espera este job.
\end{lstlisting}

% ---------------------------------------------------------------------------
\section{Figura~\ref{fig:testes} --- Estratégia de testes}

\begin{lstlisting}[language=mermaid]
flowchart TB
  v["PONTA A PONTA - Playwright, navegador real<br/>scripts/e2e/, uma pasta por PAPEL: publico, estudante, professor,<br/>gestor_escola, gestor_contrato, admin; e o probe-pointerlock"]
  c["CONTRATO EXTERNO - o que quebra em producao sem aviso<br/>telemetria-push (guard do Pub/Sub), vertex-provider (cadeia ADC e<br/>formato da requisicao), agente-professor-cliente (stream do Agent Engine)"]
  i["INTEGRACAO COM POSTGRES - o que nao da para afirmar sem banco<br/>o *-postgres.spec.ts de cada modulo, tenancy, matriculas-tenancy,<br/>planos-tenancy, motor-concorrencia, geracao-paginacao"]
  u["REGRA PURA - Vitest / Jest / pytest, milissegundos, sem banco<br/>o *-regras.spec.ts de cada modulo, motor-pesos, dossie-relatorio,<br/>bridge.test do SDK, lib/ do hub e o pytest do agente de IA do professor"]
  v --- c --- i --- u

  %% Piramide: a base (regra pura) e a mais larga, de proposito - regras.ts
  %% puro em todo modulo de dominio. E o TOPO nao e "E2E de fluxo": e
  %% verificacao de que a INTERFACE apareceu. Icone virando texto, iframe
  %% vazio, cor inexistente, fonte removida pelo Next - tudo INVISIVEL para o
  %% compilador.
\end{lstlisting}

% ---------------------------------------------------------------------------
\section{Figura~\ref{fig:cadeia-visibilidade} --- Cadeia de uma requisição e o que cada elo revela}

\begin{lstlisting}[language=mermaid]
flowchart LR
  nav[Navegador - iframe da exp] --> lb[LB + CDN] --> hub[Cloud Run -frontend] --> api[Cloud Run -api]
  api --> sql[(Cloud SQL)]
  api <--> ps[Pub/Sub]
  api --> vx[Vertex AI]

  nav -.- v0["ve-se: nada<br/>cego: crash do bundle, ponte sem hub:init, token nao renovado, lote perdido"]
  lb  -.- v1["ve-se: contagem e latencia por host<br/>cego: quem foi afetado (sem correlacao)"]
  hub -.- v2["ve-se: log de requisicao, instancias, CPU/mem<br/>cego: null do apiDoServidor, redirect em laco, erro de render"]
  api -.- v3["ve-se: idem + /readyz (se alguem abrir)<br/>cego: causa do 500, latencia da IA, descarte de telemetria, fila em erro"]
  sql -.- v4["SQL e Pub/Sub - ve-se: conexoes, CPU, disco, backlog<br/>cego: query lenta por rota; dead letter (nao existe)"]
  ps -.- v4

  %% Nenhum elo propaga identificador de correlacao.
  %% Retrato de 1o de setembro de 2026 (o ponto de partida), nao o estado
  %% corrente: por isso o Agent Engine (agente de IA do professor) nao aparece.
\end{lstlisting}

% ---------------------------------------------------------------------------
\section{Figura~\ref{fig:seq-erro-cliente} --- Erro de cliente entra pela API}

\begin{lstlisting}[language=mermaid]
sequenceDiagram
  participant E as Experiencia (@edu/sdk)
  participant H as Hub (ExpFrame)
  participant A as API (/observabilidade)
  participant L as Cloud Logging / Error Reporting

  E->>E: window.onerror / unhandledrejection
  Note over E: dedupe por hash, max. 5 por sessao, 8 KB
  E->>A: POST /observabilidade/erros-cliente (Bearer)
  Note over A: Zod: expId, versao, mensagem, stack, url
  A->>A: redaction + escola_id do token
  A->>L: stdout JSON, severity ERROR, @type ReportedErrorEvent
  L->>L: agrupa por assinatura; "novo na versao X"
  A-->>E: 204 (ou 429: cota da sessao)
  E->>H: exp:erro (opcional, targetOrigin explicito)
  Note over H: mostra "algo deu errado" e oferece recarregar
\end{lstlisting}

% ---------------------------------------------------------------------------
\section{Figura~\ref{fig:matriz-riscos} --- Matriz de riscos}

\begin{lstlisting}[language=mermaid]
quadrantChart
  title Matriz de riscos - probabilidade x impacto
  x-axis "Probabilidade baixa" --> "Probabilidade alta"
  y-axis "Impacto baixo" --> "Impacto alto"
  quadrant-1 "CRITICO - resolver antes de nov/2026"
  quadrant-2 "Impacto alto, prob. baixa"
  quadrant-3 "Monitorar"
  quadrant-4 "Prob. alta, impacto baixo"

  %% Fonte: a tabela de riscos do cap. 19. Certa conta como alta (x=0.94);
  %% possivel, como media (x=0.52).
  "R-01 carga nao verificada": [0.45, 0.91]
  "R-02 retencao indefinida": [0.82, 0.91]
  "R-03 localidade de IA": [0.94, 0.56]
  "R-04 fator humano unico": [0.94, 0.91]
  "R-05 sem alerting": [0.82, 0.56]
  "R-06 sem revisao externa": [0.82, 0.88]
  "R-07 acessibilidade so por maquina": [0.45, 0.56]
  "R-08 particao sem job": [0.15, 0.56]
  "R-09 midia manual": [0.82, 0.21]
  "R-10 orfa sem retomada": [0.45, 0.21]
  "R-11 dependencia do DS": [0.45, 0.18]
  "R-12 catalogo de modelos": [0.82, 0.53]
  "R-13 WSO2 nao e o emissor": [0.45, 0.15]
  "R-14 virada do ano letivo": [0.82, 0.85]
  "R-15 senha compartilhada": [0.45, 0.53]
  "R-16 replay sobre menor": [0.82, 0.82]
  "R-17 latencia de IA": [0.45, 0.88]
  "R-18 custo de IA exp02em1b": [0.94, 0.53]
  "R-19 guia da exp02em1b": [0.94, 0.88]
  "R-21 texto livre de menor": [0.94, 0.50]
  "R-23 turma sem missao liberada": [0.82, 0.50]
  "R-24 alavanca de incidente": [0.52, 0.91]
  "R-25 versao fixada divergindo": [0.52, 0.56]
  "R-26 historico sem prazo": [0.94, 0.21]
  "R-27 quota do Agent Engine": [0.82, 0.79]
  "R-28 ADK muda a cada 2 semanas": [0.82, 0.47]
  "R-29 stream truncado": [0.45, 0.50]
  "R-30 residencia do agente": [0.94, 0.47]
  "R-31 expire_time sem maximo": [0.45, 0.47]
  "R-32 data store nao avaliado": [0.45, 0.85]
  "R-33 Ajuda depende de vendor": [0.45, 0.44]
  "R-34 endpoint publico": [0.15, 0.53]
  "R-35 custo por minuto": [0.45, 0.56]
  "R-36 avatar mudo no erro": [0.45, 0.53]
  "R-37 base da Ajuda cresce": [0.82, 0.44]
  "R-38 painel nao e a fatura": [0.82, 0.56]
\end{lstlisting}

% ---------------------------------------------------------------------------
\section{Figura~\ref{fig:cronograma-pendencias} --- Cronograma das pendências}

\begin{lstlisting}[language=mermaid]
gantt
  title Cronograma de resolucao das pendencias (cor = prioridade)
  dateFormat YYYY-MM-DD
  axisFormat %b
  section Alta
  PEND-01 emissor para o WSO2 IS     :a1, 2026-10-01, 2026-10-31
  PEND-02 e2e do gestor regional     :a2, 2026-10-01, 2026-10-31
  PEND-05 RLS nas migracoes          :a5, 2026-09-01, 2026-09-30
  section Media
  PEND-03 alarmes de quota de IA     :m3, 2026-11-01, 2026-11-30
  PEND-04 OpenAPI no CI              :m4, 2026-09-01, 2026-09-30
  PEND-07 DLQ da telemetria          :m7, 2026-10-01, 2026-10-31
  section Baixa
  PEND-06 componentes vs. DS (continua, sem janela) :b6, 2026-09-01, 2026-12-31
  PEND-08 bundles 08 a 69            :b8, 2026-11-01, 2026-11-30
  section Marco
  Entrada em producao                :milestone, 2026-11-01, 0d
\end{lstlisting}
